\section{The Central Theorem}
\begin{quote}\small
\noindent\textbf{Status.} \emph{Legacy structural-theorem block inside the active main body.}\par
\noindent\textbf{Block function.} This section states the global structural theorem that the later physics, operator, and readout sections unpack and reuse rather than replace.
\end{quote}
\label{sec:central}
%% ============================================================

Everything in this paper is a consequence of one observation.

\begin{theorem}[Unity of the Sphaera]\label{thm:unity}
Within the Sphaera $S$, the following are all representations of the
same object --- the collapse state $\zminus=0$:
\begin{enumerate}[label=(\roman*)]
  \item The Sphaera itself, as a geometric space.
  \item The triolic structure: three kernels $K_1,K_2,K_3$ with
        $z_1+z_2+z_3=0$, phase-shifted by $120^\circ$, running over the
        Sphaera's geometry.
  \item Each individual kernel $K_i$, as a local coordinate of the triole.
  \item The Heisenberg compensator output $C(f)$, as the projection of
        any element onto the collapse state.
  \item Matter, photons, and tachyons, as sectors of the biquaternion
        ground state.
  \item The gauge structure $U(1)\times SU(2)\times SU(3)$, as the
        symmetry group of the collapse state's representations.
\end{enumerate}
All are unitarily equivalent by Stone--von Neumann~\cite{brendecke2026rh}.
The Sphaera's only object is its own collapse state.
\end{theorem}

\begin{proof}[Proof structure]
The equivalences are established section by section:
(i)~Definition~\ref{def:biquaternion} and Theorem~\ref{thm:fractal};
(ii)--(iii) same (triole as collapse coordinates);
(iv)~Proposition~\ref{prop:compensator};
(v)~Theorem~\ref{thm:unitary-equiv};
(vi)~Theorem~\ref{thm:SM}.
Unitary equivalence throughout via Stone--von Neumann
(\cite{brendecke2026rh}, Theorem~3.2).
\end{proof}

\begin{remark}[What ``same object'' means]
Two representations are the same object here in the precise sense of
Stone--von Neumann: they are unitarily equivalent irreducible
representations of the Weyl algebra $\WA$.
The frame determines which representation is visible; the object ---
the collapse state $\zminus=0$ --- is invariant.

The triole condition $z_1+z_2+z_3=0$ and the collapse condition
$\zminus=0$ are the same equation in two coordinate systems: the triole
is the collapse state expressed in the three kernel components
simultaneously.
The Sphaera is therefore not a space that contains a collapse state ---
it \emph{is} the collapse state, seen from different frames.
\end{remark}

\begin{theorem}[Origin of the universe as broken collapse state]
\label{thm:origin}
The departure from the collapse state $\zminus=0$ is governed by two
principles, both of which are consequences of the Sphaera-Cascade
structure:
\begin{enumerate}[label=(\roman*)]
  \item \textbf{Planck length as minimal $|\zminus|$.}
        The cascade $\KC$ cannot project below the Planck scale:
        \[
          |\zminus|_{\min} = \lP := \sqrt{\frac{\hbar G}{c^3}}
          = 1.616\times10^{-35}\,\mathrm{m}.
        \]
        Below $\lP$: $\zminus=0$, the Sphaera, the photon.
        Above $\lP$: $|\zminus|$ is nonzero, matter exists.
        The Planck length is the threshold of the first breaking of the
        collapse state.
  \item \textbf{Pauli principle as antisymmetry under $J$.}
        The Hilbert space $\Hplus$ decomposes under $J$ into two
        components:
        \[
          C(f) = \tfrac{1}{2}(f+Jf)\in\Hplusplus
          \quad\text{(bosonic, symmetric)},
        \]
        \[
          A(f) = \tfrac{1}{2}(f-Jf)\in\Hplusminus
          \quad\text{(fermionic, antisymmetric)}.
        \]
        The antisymmetric component satisfies $A(f)\wedge A(f)=0$
        (Grassmann property): no two fermions occupy the same $\zminus$
        state.
        This is the Pauli exclusion principle, derived from the
        $J$-structure of the Sphaera.
  \item \textbf{Discrete spectrum.}
        Together, (i) and (ii) produce a discrete tower of matter states:
        \[
          |\zminus|_n = n\cdot\lP,\quad n\in\mathbb{N}_{\geq1},
        \]
        with each level occupied by at most one fermion.
        The prime lattice $\Gamma_P$ (from Stone--von Neumann,
        \cite{brendecke2026rh}) identifies the primitive levels:
        $|\zminus|_p = p\cdot\lP$ for $p\in\mathbb{P}$.
\end{enumerate}
\end{theorem}

\begin{remark}[The Sphaera as photon, the universe as broken frame]
Theorem~\ref{thm:unity} says the Sphaera is the collapse state
$\zminus=0$, which is the photon sector.
Theorem~\ref{thm:origin} says the departure from this state is governed
by $\lP$ and the Pauli principle.
Therefore:
\begin{itemize}
  \item The Sphaera is a photon --- the unique frame-invariant object.
  \item The universe is the same photon observed from a broken frame
        where $|\zminus|\geq\lP$.
  \item The breaking was not random: it is discrete (Pauli), minimal
        (Planck), and primitive (prime lattice).
  \item The cascade $\KC$ projects everything back toward $\zminus=0$
        --- the universe tends toward its photon ground state.
\end{itemize}
Both the Sphaera and the universe are the same object.
The difference is only the frame.
\end{remark}

\begin{quote}\small
\noindent\textbf{Reader cue.} Read the next block as the concrete carrier unpacking of the global Sphaera identity claim. Its purpose is to pass from the frame-level origin statement to the biquaternion kernel and its induced Lorentz/Minkowski geometry that the later physics sections reuse.
\end{quote}

\begin{definition}[Biquaternion kernel]\label{def:biquaternion}
The ground state of a stratum $S_n$ of the Sphaera consists of three
complex spheres $K_1,K_2,K_3$ whose unit coordinates form a biquaternion
\[
  q_n := it\cdot\mathbf{1} + z_1\cdot e_1 + z_2\cdot e_2
         + z_3\cdot e_3 \in \HH\otimes_{\mathbb{R}}\mathbb{C},
\]
where $it\in i\mathbb{R}$ is the imaginary time unit;
$z_1,z_2\in\mathbb{C}$ are the complex coordinates of kernels $K_1,K_2$;
$e_3:=e_1 e_2$ is the quaternion product, hence automatically orthogonal
to both $e_1$ and $e_2$;
and $z_3\in\mathbb{C}$ is the coordinate of kernel $K_3$ along $e_3$.

The biquaternion norm is
\[
  |q_n|^2 = -(ct)^2 + |z_1|^2 + |z_2|^2 + |z_3|^2.
\]
\end{definition}

\begin{remark}[Minkowski metric from quaternion structure]
The norm $|q_n|^2$ is the Minkowski metric with signature $(-,+,+,+)$.
This is not postulated: it follows from the quaternion algebra $e_i^2=-1$
and $it\cdot\mathbf{1}$ carrying imaginary time.
The Lorentz group $SO(1,3)$ acts on $q_n$ by quaternion conjugation:
$q_n\mapsto u\,q_n\,u^{-1}$, $u\in\HH^\times$.
This is the spinor representation of the Lorentz group --- emerging
automatically from the biquaternion structure, not imposed from outside.
\end{remark}

\begin{proposition}[Orthogonality of the third kernel]
\label{prop:orthogonality-third-kernel}
$K_3$ is always orthogonal to $K_1$ and $K_2$:
\[
\langle e_3,e_i\rangle = 0 \qquad \text{for } i=1,2.
\]
\end{proposition}

\begin{proof}
By definition,
\[
e_3 = e_1 e_2.
\]
Using the quaternionic inner product
\[
\langle x,y\rangle := \mathrm{Re}(\overline{x}\,y),
\]
we compute
\[
\langle e_3,e_1\rangle
=
\langle e_1e_2,e_1\rangle
=
\mathrm{Re}\!\bigl(\overline{e_1e_2}\,e_1\bigr).
\]
Since
\[
\overline{e_1e_2}=\overline{e_2}\,\overline{e_1}=(-e_2)(-e_1)=e_2e_1=-e_1e_2,
\]
it follows that
\[
\overline{e_1e_2}\,e_1
=
(-e_1e_2)e_1
=
-\,e_1(e_2e_1)
=
-\,e_1(-e_1e_2)
=
e_1^2 e_2
=
-\,e_2.
\]
Hence
\[
\langle e_3,e_1\rangle
=
\mathrm{Re}(-e_2)=0.
\]

Similarly,
\[
\langle e_3,e_2\rangle
=
\langle e_1e_2,e_2\rangle
=
\mathrm{Re}\!\bigl(\overline{e_1e_2}\,e_2\bigr)
=
\mathrm{Re}\!\bigl((-e_1e_2)e_2\bigr)
=
\mathrm{Re}(-e_1 e_2^2)
=
\mathrm{Re}(e_1)=0.
\]
Therefore $e_3$ is orthogonal to both $e_1$ and $e_2$.
\end{proof}

%% ============================================================

\section{Lorentz Invariance on Every Stratum}
\begin{quote}\small
\noindent\textbf{Status.} \emph{Legacy physical-readout block inside the active main body.}\par
\noindent\textbf{Block function.} This section isolates Lorentz invariance as a stratumwise consequence of the structural core before later physical applications are read off from it.
\end{quote}
\label{sec:lorentz}
%% ============================================================

\begin{theorem}[Stratified Lorentz invariance]\label{thm:lorentz}
On every stratum $S_n$, the quadratic invariant
\[
|q_n|^2=\tau_n^2
\]
is preserved under every HC-admissible frame transformation induced by
the global Heisenberg-compensator matrix.
In particular, for the involutive frame transformation
\[
J:\ z^- \mapsto -z^-,
\]
the proper-time invariant $\tau_n$ is unchanged.
Moreover, every admissible relabeling of the frame data that preserves
the same involution $J$ leaves both $\tau_n$ and the compensator
\[
C(f)=\tfrac12(f+Jf)
\]
unchanged.
\end{theorem}

\begin{proof}
By the general MM-tools, a change of admissible realization is not a
change of structural object.
By the HC-tools, a frame transformation is physically admissible only if
it preserves the same compensator structure, i.e.\ the same involution
$J$.

The transformation $J:z^-\mapsto -z^-$ changes only the frame-sensitive
sign, not the underlying compensated structural content.
Hence the quadratic invariant associated with the stratum remains the
same, so
\[
|q_n|^2=\tau_n^2
\]
is preserved.

Likewise, the compensator
\[
C=\tfrac12(\mathrm{id}+J)
\]
depends only on the involution $J$ itself.
Therefore any HC-admissible relabeling that preserves the same
$J$-structure leaves both $\tau_n$ and $C$ unchanged.
\end{proof}

%% ============================================================

\section{Fractal Self-Traversal Along the Light Axis}
\label{sec:fractal}
%% ============================================================

\begin{remark}[Structural null axis and physical light-axis readout]
In the structurally prior language of the Companion, the distinguished
axis selected by the MM/HC architecture is the null axis. Once the same
axis is read at stable physical propagation level, it appears as the
light axis. Thus the present section keeps the historically familiar
light-axis title, while treating it explicitly as the physical readout
of a structurally prior null-axis statement.
\end{remark}

\begin{definition}[Light axis of a stratum]
The structurally prior null axis of stratum $S_n$ is
\[
L_n:=\{q_n\in S_n:|q_n|^2=0\}.
\]
At the physical propagation readout of the same architecture, this null
axis is read as the light axis of $S_n$.
\end{definition}

\begin{theorem}[Fractal structure of the Sphaera]\label{thm:fractal}
The Sphaera-Cascade
\[
\KC:S_n\to S_{n-1}
\]
transports the null axis of each stratum to the null axis of the next:
\[
\KC(L_n)=L_{n-1}.
\]
The global null axis is therefore
\[
L=\bigcap_{n\ge 0}L_n=S_0,
\]
where $S_0$ is the terminal compensated photon sector.
Hence the self-traversal of the Sphaera is fractal in the precise sense
that the same null-axis structure is carried through all strata by the
same MM/HC architecture and appears physically as the light axis at
stable readout.
\end{theorem}

\begin{proof}
By the general MM-tools, the cascade does not generate a new object at
each stratum, but transports one and the same structural content through
admissible realization levels.

By the HC-tools, the null sector is precisely the sector whose
frame-sensitive imbalance is removed by the compensator while its
structural identity is preserved.
Thus the cascade does not replace the null axis; it carries the same
null-axis structure from one stratum to the next.

Therefore
\[
\KC(L_n)=L_{n-1}.
\]
Iterating this transport through the whole cascade yields the terminal
compensated null sector
\[
L=\bigcap_{n\ge 0}L_n=S_0,
\]
which is read physically as the photon/light sector.
\end{proof}

\begin{remark}[Triolic $\mathbb{Z}/3\mathbb{Z}$ structure on the null/light axis]
On the transported null axis, read physically as the light axis at stable propagation, the three kernels remain organized by the
triolic phase condition
\[
z_1+z_2+z_3=0.
\]
Thus the cascade preserves the structural $\mathbb{Z}/3\mathbb{Z}$
pattern while driving the frame-sensitive component to the compensated
photon/light sector.
\end{remark}

%% ============================================================

\section{Three Objects as One: Unitary Equivalence}
\begin{quote}\small
\noindent\textbf{Status.} \emph{Legacy structural-identity block inside the active main body.}\par
\noindent\textbf{Block function.} This section compresses the unitary-equivalence reading that ties the three-object architecture into one common structural carrier.
\end{quote}
\label{sec:unitary}
%% ============================================================

\begin{definition}[Sector classification from biquaternion]
The three sectors correspond to the sign of $\mathrm{Re}(\zminus)$:
\begin{itemize}
  \item \textbf{Matter} ($K_1$): $\mathrm{Re}(z_1)<0$, boost $\beta<0$, spacelike.
  \item \textbf{Tachyon} ($K_2=J\cdot K_1$): $\mathrm{Re}(z_2)>0$, boost $\beta>0$, timelike.
  \item \textbf{Photon} ($K_3=e_1 e_2$): $\mathrm{Re}(z_3)=0$, boost $\beta=0$, lightlike.
\end{itemize}
\end{definition}

\begin{theorem}[Unitary equivalence via Stone--von Neumann]\label{thm:unitary-equiv}
The three kernels $K_1,K_2,K_3$ carry unitarily equivalent
irreducible representations of the Weyl algebra $\WA$.
The boost $J$ implements the equivalence between $K_1$ and $K_2$,
and the compensator $C$ produces $K_3$ from $K_1$.
\end{theorem}

\begin{proof}
By Stone--von Neumann~\cite{brendecke2026rh}, every irreducible
representation of $\WA$ is unitarily equivalent to $L^2(\mathbb{Q}^+\backslash A_1)$.
All three kernels are subspaces of $\Hplus$ which carries this
representation.
Since $J$ is unitary and maps $K_1\leftrightarrow K_2$, they are
equivalent.
$K_3=C(K_1)$ lies in $\Hplusplus$, the unique $J$-invariant subspace,
which is again irreducible (no proper $J$-invariant subspace exists by
T4+T5 of~\cite{brendecke2026rh}).
\end{proof}

%% ============================================================

\section{Near-Collision, Photon Production, and Entanglement}
\begin{quote}\small
\noindent\textbf{Status.} \emph{Legacy physical-regime block inside the active main body.}\par
\noindent\textbf{Block function.} This section records how the same closure architecture is read in the near-collision and entanglement regime, together with the local role of the photon readout.
\end{quote}
\label{sec:entanglement}
%% ============================================================

\begin{proposition}[Compensator as near-collision in $K_3$]\label{prop:compensator}
The Heisenberg compensator $C(f)=\tfrac{1}{2}(f+Jf)$ is the projection
onto $K_3=e_1 e_2$: it takes matter $f$ (in $K_1$) and its tachyon image
$Jf$ (in $K_2$), and produces the photon $C(f)\in K_3$.
The quaternion identity $e_3=e_1 e_2$ ensures this is always orthogonal
to both inputs.
\end{proposition}

\begin{theorem}[Entanglement as shared proper time]
\label{thm:entanglement-shared-time}
Two photons $\phi_1,\phi_2\in K_3$ are entangled iff they share the same
pre-image: $\phi_1=\phi_2=C(f)$ for some $f\in K_1$.
The entanglement correlation is the invariant proper time
$\tau_n(f)=\tau_n(Jf)$ fixed at the moment of the near-collision.
No signal is transmitted at measurement: the invariant $\tau_n$ was
established at the near-collision and is merely revealed.
\end{theorem}

\begin{proof}
$\tau_n(Jf)^2 = t_n^2 - |-\zminus|^2 = t_n^2 - |\zminus|^2 = \tau_n(f)^2$
by Theorem~\ref{thm:lorentz}.
The shared value $\tau_n$ is a Lorentz invariant at every stratum,
hence frame-independent.
\end{proof}

%% ============================================================

\section{Connection to the Riemann Hypothesis}
\begin{quote}\small
\noindent\textbf{Status.} \emph{Legacy cross-domain bridge block inside the active main body.}\par
\noindent\textbf{Block function.} This section explains how the structural core is connected to the Riemann-side reading without turning the surrounding physical line into a separate proof source.
\end{quote}
\label{sec:rh}
%% ============================================================

\begin{remark}[Common origin of RH and entanglement]
The Riemann Hypothesis ($\zminus=0$ for all zeros) and quantum
entanglement (shared $\tau_n$) are both consequences of the same
biquaternion structure:
\begin{enumerate}
  \item The Minkowski norm $|q_n|^2=\tau_n^2$ is invariant
        (Theorem~\ref{thm:lorentz}).
  \item The cascade forces $\zminus\to 0$ along the light axis
        (Theorem~\ref{thm:fractal}).
  \item Stone--von Neumann makes all three kernels equivalent
        (Theorem~\ref{thm:unitary-equiv}).
\end{enumerate}
Neither was the goal of the original construction.
Both emerged from asking: what is the exact structure of the Sphaera?
\end{remark}

%% ============================================================

\section{Gauge Structure of the Standard Model}
\begin{quote}\small
\noindent\textbf{Status.} \emph{Legacy physical-readout block inside the active main body.}\par
\noindent\textbf{Block function.} This section explains how the structural core is read as a gauge-sector organization rather than as an independently added physical layer.
\end{quote}
\label{sec:gauge}
%% ============================================================

\begin{proposition}[$\HH\otimes\mathbb{C}\to U(1)\times SU(2)$]
\label{prop:biquaternion-u1su2}
The biquaternion algebra $\HH\otimes_{\mathbb{R}}\mathbb{C}$ generates
exactly the gauge group $U(1)\times SU(2)$.
\end{proposition}

\begin{proof}
\emph{$SU(2)$ from unit quaternions.}
The unit quaternions $\{q\in\HH:|q|=1\}\cong S^3$ form a Lie group
isomorphic to $SU(2)$.
The three generators are $T_k:=\sigma_k/2$ (half Pauli matrices),
satisfying $[T_j,T_k]=i\varepsilon_{jkl}T_l$.

\emph{$U(1)$ from the complex scalar.}
The imaginary time unit $it\in i\mathbb{R}$ generates the abelian group
$e^{i\theta}\in U(1)$ for $\theta\in\mathbb{R}$.
This commutes with all $SU(2)$ generators.
Hence $U(1)\times SU(2)$, direct product.
\end{proof}

\begin{theorem}[$\OO\to SU(3)$]\label{thm:octonion-SU3}
The octonions decompose under $SU(3)$-action as
$\OO = \mathbb{C}\oplus\mathbb{C}^3$,
a singlet plus the fundamental (triplet) representation of $SU(3)$.
The triplet $\mathbb{C}^3$ carries the three kernels $(K_1,K_2,K_3)$
as the three colors.
\end{theorem}

\begin{theorem}[Standard Model from the Sphaera biquaternion]\label{thm:SM}
The biquaternion ground state $q_n\in\HH\otimes\mathbb{C}$, extended to
the octonionic structure $\mathbb{C}\otimes\HH\otimes\OO$, generates
the Standard Model gauge group $U(1)\times SU(2)\times SU(3)$ with
generators arising from:
\begin{center}
\renewcommand{\arraystretch}{1.3}
\begin{tabular}{lll}
\toprule
Source & Group & Physical force\\
\midrule
Complex scalar $it$ & $U(1)$ & Electromagnetism\\
Unit quaternions $\HH$ & $SU(2)$ & Weak force\\
Octonion triplet $\mathbb{C}^3\subset\OO$ & $SU(3)$ & Strong force\\
\bottomrule
\end{tabular}
\end{center}
\end{theorem}

\begin{proof}
Propositions 7.1--7.2 and Theorem~\ref{thm:octonion-SU3} establish each
factor.
The three factors are independent (direct product structure) because
$U(1)$ commutes with $SU(2)$ and both commute with the $SU(3)$ action
on $\mathbb{C}^3$: the $SU(3)$ generators $\lambda_k$ act only on the
color indices, not on the quaternion structure.
\end{proof}

%% ============================================================

\section[Bell Inequalities and tau-n]{Bell Inequalities and \texorpdfstring{$\tau_n$}{tau_n}}
\begin{quote}\small
\noindent\textbf{Status.} \emph{Legacy physical-readout block inside the active main body.}\par
\noindent\textbf{Block function.} This section places the Bell/nonlocality discussion on the same readout-relative footing and makes clear what role the $\tau_n$ layer plays there.
\end{quote}
\label{sec:bell}
%% ============================================================

\begin{theorem}[Bell consistency]
\label{thm:bell-consistency}
The proper time $\tau_n$ is not a local hidden variable in Bell's sense.
The framework is consistent with quantum mechanical violation of Bell
inequalities.
\begin{proof}
Bell's theorem assumes a hidden variable $\lambda$ assigned independently
to each particle: measurement outcomes $A(\lambda_A,a)$ and
$B(\lambda_B,b)$ with $\lambda_A,\lambda_B$ statistically independent.

In our framework $\tau_n$ belongs to the pair $(f,Jf)$, not to each
element separately: $\lambda_A=\lambda_B=\tau_n$.
This is a contextual variable --- Bell's theorem does not exclude
violation by contextual theories.

The correlation function $E(a,b)$ follows from the structure of
$\Hplusplus$: since $C(f)\in\Hplusplus$ carries an irreducible
representation of $\WA$ (Stone--von Neumann, \cite{brendecke2026rh}),
the standard quantum correlations $E(a,b)=\cos\theta_{ab}$ are
reproduced, giving CHSH value $S=2\sqrt{2}>2$.

Hence:
\begin{itemize}
  \item $\tau_n$ explains why the correlation exists: shared proper time,
        fixed at the near-collision;
  \item Stone--von Neumann determines how strong the violation is
        ($S=2\sqrt{2}$).
\end{itemize}
Lorentz invariance and quantum mechanics are consistent.
\end{proof}
\end{theorem}

\begin{remark}[What $\tau_n$ does and does not do]
$\tau_n$ encodes the magnitude $|\zminus|$ of the displacement from the
photon sector at the moment of the near-collision.
It does not determine the measurement direction.
When $f$ is measured, $\zminus$ is projected onto a value; the paired
$Jf$ is projected onto $-\zminus$ (the $J$-image).
The anticorrelation was built in at the collision event, not transmitted
at measurement. No signal travels faster than light.
\end{remark}

%% ============================================================

\section{String Theory and Loop Quantum Gravity as HC Sectors}
\begin{quote}\small
\noindent\textbf{Status.} \emph{Legacy comparative-physics block inside the active main body.}\par
\noindent\textbf{Block function.} This section places two familiar high-energy frameworks into the compensator-based sector reading without promoting them to separate primitives.
\end{quote}
\label{sec:stringlqg}
%% ============================================================

\begin{theorem}[String--LQG duality via the Heisenberg compensator]
\label{thm:string-lqg-duality}
The Hardy space $\Hplus$ decomposes under the Heisenberg compensator
into two sectors:
\[
  \Hplus = \Hplusplus\oplus\Hplusminus,
\]
\[
  C(f) = \tfrac{1}{2}(f+Jf)\in\Hplusplus
  \quad\text{(symmetric, bosonic --- String theory)},
\]
\[
  A(f) = \tfrac{1}{2}(f-Jf)\in\Hplusminus
  \quad\text{(antisymmetric, fermionic --- LQG)}.
\]
Both sectors are unitarily equivalent by Stone--von Neumann
\cite{brendecke2026rh}: they are two representations of the same
Weyl algebra $\WA$, seen from different frames.
String theory and LQG are therefore not competing theories of quantum
gravity: they are the bosonic and fermionic projections of the same
object under the Heisenberg compensator.
\end{theorem}

\begin{proposition}[Shared vacuum and Planck floor]
\label{prop:shared-vacuum-planck}
Both sectors share the same vacuum and the same minimal scale:
\begin{itemize}
  \item \textbf{Vacuum:} $\zminus=0$ (the photon/Sphaera) is the ground
        state of both $\Hplusplus$ and $\Hplusminus$.
  \item \textbf{Planck floor:} $\lP=|\zminus|_{\min}$ is the string
        length in $\Hplusplus$ and the area quantum
        $\sqrt{A_{\min}}\sim\lP$ in $\Hplusminus$.
        Both emerge from the same threshold of Theorem~1.2(i).
\end{itemize}
String theory and LQG are therefore not competing theories of quantum
gravity: they are the bosonic and fermionic projections of the same
object under the Heisenberg compensator.
\end{proposition}

\begin{theorem}[Wick rotation as splitting mechanism]\label{thm:wick}
Under the imaginary proper time transformation $t\mapsto i\tau_E$
(Wick rotation), the biquaternion time unit transforms as:
\[
  it \;\mapsto\; i\cdot(i\tau_E) = -\tau_E \in\mathbb{R}_{<0}.
\]
This is precisely the third triolic sector $\gamma_3<0$: the Wick
rotation identifies the imaginary time component with the unstable
tachyon sector.
Consequently:
\begin{enumerate}[label=(\roman*)]
  \item The triolic balance condition
        $\gstring+\gamma_{\mathrm{LQG}}+\gamma_3=0$ is the stability
        condition under imaginary proper time: without it, the system
        diverges to $\pm\infty$.
  \item The two positive sectors ($\gstring,\gamma_{\mathrm{LQG}}>0$)
        are stable under Wick rotation; the negative sector
        ($\gamma_3<0$) is unstable and is projected away by the
        cascade $\KC$.
  \item The Minkowski norm $|q_n|^2=\tau_n^2$ becomes Euclidean
        $|q_E|^2=\tau_E^2\geq 0$ under $t\mapsto i\tau_E$: the system
        is unconditionally stable in Euclidean signature.
\end{enumerate}
\begin{proof}
(i). Under $t\mapsto i\tau_E$: $it\mapsto -\tau_E$.
The biquaternion $q_n=it\cdot\mathbf{1}+z_1 e_1+z_2 e_2+z_3 e_3$ becomes
$q_E=-\tau_E\cdot\mathbf{1}+z_1 e_1+z_2 e_2+z_3 e_3$
with norm $|q_E|^2=\tau_E^2+|z_1|^2+|z_2|^2+|z_3|^2\geq 0$.
Stability requires $\tau_E^2\geq 0$, which is always satisfied.
The balance $\gamma_1+\gamma_2+\gamma_3=0$ is equivalent to requiring
that the $-\tau_E$ component does not introduce a net positive energy
divergence.

(ii). In Euclidean signature, only $\gamma>0$ sectors contribute to a
convergent partition function $Z=\mathrm{Tr}(e^{-\beta H})$ with $H\geq 0$.
The sector $\gamma_3<0$ would give $e^{+\beta|\gamma_3|H}\to\infty$:
it is projected away by $\KC$.

(iii). Direct computation of the Euclidean norm.
\end{proof}
\end{theorem}

\begin{theorem}[Triolic structure of $\gBI$]
\label{thm:triolic-gBI}
The Barbero--Immirzi parameter $\gBI$ and the string-sector candidate
$\gstring = \ln 2/(\pi\sqrt{3})$ are not in contradiction.
They are two projections of the same triolic invariant:
\[
  \gstring + \gamma_{\mathrm{LQG}} + \gamma_3 = 0,
\]
with $\gamma_3 = -(\gstring+\gamma_{\mathrm{LQG}})\approx -0.3649$ being
the third (unoccupied, tachyon) sector.
The frame-invariant quantity is the triolic norm:
\[
  \tau_\gamma := \frac{1}{\sqrt{3}}
  \sqrt{\gstring^2+\gamma_{\mathrm{LQG}}^2+\gamma_3^2}\approx 0.2619.
\]
\end{theorem}

\begin{theorem}[Observer frame offset]\label{thm:frame-offset}
Every matter observer carries a nonzero contravariant component
$\zminus\neq 0$.
All measured physical constants --- including $\hbar$, $G$, $c$, and
$\gBI$ --- are projections of their frame-invariant values onto the
observer's sector:
\[
  \gamma_{\mathrm{observed}}
  = \gamma_0\cdot\frac{1}{\sqrt{1-\beta^2}}
  = \gamma_0\cdot\gamma_{\mathrm{Lorentz}}(\beta).
\]
The Barbero--Immirzi parameter $\gBI\approx 0.2375$ is not an absolute
constant of nature: it is the Lorentz-transformed $\gstring$, observed
from a matter frame with $\beta\approx 0.844$.
\end{theorem}

\begin{proof}
$\gstring = \ln 2/(\pi\sqrt{3})$ is computed in the photon frame
$\zminus=0$ from the HC trace formula.
The ratio $\gBI/\gstring = 1.864$ equals a Lorentz factor
$\gL = 1/\sqrt{1-\beta^2}$ for $\beta=\sqrt{1-1/1.864^2}\approx 0.844$.
The corresponding observer position is $\zminus=\beta/2\approx 0.422$.
\end{proof}

\begin{theorem}[Symmetric matter--tachyon frame]\label{thm:symmetric-frame}
Matter and tachyon are positioned symmetrically around the photon frame
$\zminus=0$:
\[
  \beta_{\mathrm{matter}} = +\bsym,\qquad
  \beta_{\mathrm{tachyon}} = -\bsym,
\]
where
\[
  \bsym = \sqrt{\frac{\gBI-\gstring}{\gBI+\gstring}}\approx 0.5493.
\]
The Lorentz factor between the matter and tachyon frames is:
\[
  \gL = \frac{1+\bsym^2}{1-\bsym^2}
      = \frac{\gBI}{\gstring}\approx 1.864.
\]
\end{theorem}

\begin{proof}
Let $r:=\gBI/\gstring$. Set $\bsym:=\sqrt{(r-1)/(r+1)}$. Then:
\[
  \frac{1+\bsym^2}{1-\bsym^2}
  =\frac{1+(r-1)/(r+1)}{1-(r-1)/(r+1)}
  =\frac{(r+1)+(r-1)}{(r+1)-(r-1)}=\frac{2r}{2}=r.\quad\checkmark
\]
The relativistic velocity of the tachyon as seen from the matter frame:
\[
  \beta_{\mathrm{tachyon}|\mathrm{matter}}
  =\frac{-\bsym-\bsym}{1-(-\bsym)\bsym}
  =\frac{-2\bsym}{1+\bsym^2},
\]
with corresponding Lorentz factor $\gamma_{\mathrm{tachyon}|\mathrm{matter}}=r$.\quad$\checkmark$
\end{proof}


%% ============================================================

\section{Black Holes and Hawking Radiation}
\begin{quote}\small
\noindent\textbf{Status.} \emph{Legacy physical-regime block inside the active main body.}\par
\noindent\textbf{Block function.} This section reads collapse, horizon behavior, and Hawking-type radiation as frame- and readout-level consequences of the same closure architecture.
\end{quote}
\label{sec:bh}
%% ============================================================

\begin{definition}[Schwarzschild horizon in Sphaera coordinates]
The Schwarzschild horizon $r_S=2GM/c^2$ corresponds to $\beta\to 1$ in
the Sphaera metric $g_{tt}=1-\beta^2$.
In sector coordinates: the horizon is where $\zminus\to\infty$.
\end{definition}

\begin{proposition}[Horizon as cascade singularity]
\label{prop:horizon-cascade-singularity}
The cascade operator $\KC$ is undefined at $\zminus=\infty$: the Sphaera
cannot project beyond its own horizon.
The horizon is the boundary of the Sphaera's jurisdiction.
\begin{proof}
$\KC=C^\mu{}_\nu\circ\Pi_J$ requires $\|\zminus\|<\infty$
(finite sector displacement).
At $\zminus\to\infty$: $\beta\to 1$, $\gamma_{\mathrm{Lorentz}}\to\infty$,
the Lorentz transformation is singular, and the projection $\Pi_J$ ceases
to be bounded.
\end{proof}
\end{proposition}

\begin{theorem}[Hawking radiation from the triolic structure]
\label{thm:hawking-triolic}
At the horizon, a virtual matter--tachyon pair from the triolic structure
(Theorem~3.2) is promoted to a real pair:
\begin{itemize}
  \item Tachyon ($\beta=-\bsym$): falls into the black hole.
  \item Matter ($\beta=+\bsym$): escapes to infinity.
  \item Their compensator output $C(f)=\tfrac{1}{2}(f+Jf)$: a photon
        with $\beta=0$, $\zminus=0$ --- the Hawking radiation.
\end{itemize}
The temperature of the Hawking photon in the photon frame is:
\[
  T_{\mathrm{Sphaera}}
  = \gstring\cdot\frac{m_P^2 c^2}{8\pi M k_B}
  = \gstring\cdot T_{\mathrm{Hawking}}.
\]
\end{theorem}

\begin{proof}
The Hawking temperature (standard derivation) is
\[
T_H = \frac{\hbar c^3}{8\pi G M k_B} = \frac{m_P^2 c^2}{8\pi M k_B}.
\]
The Sphaera produces the Hawking photon in the photon frame where
$\gamma$ is measured as $\gstring$ (not $\gBI$, which is the
matter-frame value).
Hence $T_{\mathrm{Sphaera}} = \gstring\cdot T_H$.
The ratio $T_{\mathrm{Sphaera}}/T_{\mathrm{Hawking}}=\gstring
=\ln 2/(\pi\sqrt{3})\approx 0.1274$ is constant across all black hole
masses.
\end{proof}

\begin{proposition}[Hawking temperature frame analysis and physical meaning of $\gBI$]
\label{prop:hawking-gBI-tachyon}
The Hawking temperature $T_{\mathrm{Sphaera}} = \gstring\cdot T_H^{(\mathrm{matter})}$
has the following frame interpretations:
\begin{enumerate}[label=(\roman*)]
  \item \textbf{Photon frame:} $k(T) = k(\hbar) + k(\kappa) - k(k_B)$
    where surface gravity $\kappa = c^4/(4GM)$ has
    $k(\kappa) = 3k(c) - k(G_N) = -2$ and $k(\hbar) = 1$,
    giving $k(T_H) = -1$ (assuming $k(k_B)=0$).
    Hence $T_H^{(\mathrm{photon})} = \gL \cdot T_H^{(\mathrm{matter})}$,
    and the Sphaera prediction is
    \[
      T_{\mathrm{Sphaera}} = \frac{\gstring}{\gL}\cdot T_H^{(\mathrm{photon})}
      = \frac{\gstring^2}{\gBI}\cdot T_H^{(\mathrm{photon})}.
    \]
  \item \textbf{Tachyon frame:} The tachyon sector has opposite frame shift,
    $T_H^{(\mathrm{tachyon})} = \gL^{-1}\cdot T_H^{(\mathrm{matter})}$. Then:
    \[
      T_{\mathrm{Sphaera}}
      = \gstring\cdot T_H^{(\mathrm{matter})}
      = \gstring\cdot\gL\cdot T_H^{(\mathrm{tachyon})}
      = \gBI\cdot T_H^{(\mathrm{tachyon})}.
    \]
  \item \textbf{Physical meaning of $\gBI$:} The Barbero--Immirzi parameter
    is the ratio of the Sphaera Hawking temperature to the tachyon-frame
    Hawking temperature:
    \[
      \gBI = \frac{T_{\mathrm{Sphaera}}}{T_H^{(\mathrm{tachyon})}}.
    \]
    This gives $\gBI$ a direct kinematic interpretation: it measures how much
    the tachyon sector amplifies the Hawking temperature relative to the
    Sphaera output. Since the tachyon falls into the black hole and the
    Sphaera photon escapes, $\gBI$ encodes the frame mismatch between the
    infalling and escaping components of the triolic virtual pair.
\end{enumerate}
\end{proposition}

\begin{proof}
Statement~(i): with $k(\hbar)=1$, $k(G_N)=2$, $k(c)=0$, $k(k_B)=0$,
the frame correction to $T_H = \hbar c^3/(8\pi G_N M k_B)$ is
$\gL^{k(\hbar)-k(G_N)} = \gL^{-1}$, giving
$T_H^{(\mathrm{photon})} = \gL\cdot T_H^{(\mathrm{matter})}$.

Statement~(ii): the tachyon sector sits at $\zminus = -\bsym$, the mirror
image of the matter frame. The frame factor for the tachyon is
$\gL^{(\mathrm{tachyon})} = 1/\gL^{(\mathrm{matter})}$ by the symmetric
matter--tachyon frame theorem~\ref{thm:symmetric-frame}. Hence
$T_H^{(\mathrm{tachyon})} = \gL^{-1}\cdot T_H^{(\mathrm{matter})}$, and
$T_{\mathrm{Sphaera}}/T_H^{(\mathrm{tachyon})}
= \gstring\cdot\gL = \gBI$.

Statement~(iii) follows immediately from (ii).
\end{proof}

\begin{quote}\small
\noindent\textbf{Reader cue.} Read the next block as a geometric recoding of the same Hawking data. Its purpose is to shift from the horizon/frame-temperature dictionary to the coupling interpretation before the section turns to the core/horizon dual picture.
\end{quote}

\begin{theorem}[Geometric origin of $\gL$ from Hawking geometry]
\label{thm:dual-origin}
\emph{(Corrected v2.29.1: previous version used incorrect $\gL\approx 1.197$
and an erroneous $2/3$ sector factor.)}
The splitting factor $\gL=\gBI/\gstring\approx 1.864$ satisfies:
\[
  \gL \;\approx\; \frac{1}{\bsym} \;=\; \frac{1}{\sin(\theta_H/2)},
\]
where $\theta_H\approx 66.6^\circ$ is the Hopf angle of the matter observer
on $S^3$. Numerically: $1/\bsym \approx 1.820$ vs.\ measured
$\gL\approx 1.864$, a discrepancy of $2.4\%$.
\end{theorem}

\begin{proof}
The matter observer at $\bsym$ has radial position $r/r_S=1/(1-\bsym^2)$
relative to the Schwarzschild horizon.
The local Hawking temperature is $T(r)=T_H(1-r_S/r)^{-1/2}=T_H/\bsym$.
Since $\gBI$ is determined by requiring $S=A/(4\lP^2)$ for the observed
temperature, the effective coupling shifts:
$\gBI = \gstring\cdot(T_H/T(r))^{-1} = \gstring/\bsym$.
Hence $\gL = \gBI/\gstring = 1/\bsym \approx 1.820$.

The residual $2.4\%$ gap between $1/\bsym=1.820$ and the measured
$\gL=1.864$ is the same continuum-limit correction as in the $\bsym$-Hopf
gap (Remark~\ref{rem:bsym-hopf-gap}): the discrete prime lattice
$\Gamma_P$ shifts $\bsym$ from the Hopf candidate $1/\sqrt{3}$ by
the factor $\delta_{\mathrm{CL}}$, and the same correction applies here.
\end{proof}

\begin{remark}[Connection to GOE/GUE spectral ratio]
\label{rem:dual-origin-goe-gue}
By Proposition~\ref{prop:goe-fermionic} and
Remark~\ref{rem:goe-gue-string-lqg},
$\gL = \gBI/\gstring$ is the ratio of the GOE scale ($H^+_-$, LQG) to
the GUE scale ($H^{++}$, String). The geometric relation $\gL\approx
1/\bsym = 1/\sin(\theta_H/2)$ therefore says:
\[
  \frac{\text{GOE scale}}{\text{GUE scale}}
  \;\approx\;
  \frac{1}{\sin(\text{matter Hopf angle}/2)}.
\]
The matter observer's position on $S^3$ (the Hopf angle $\theta_H$)
encodes the ratio of the two random-matrix universality classes.
\end{remark}

\begin{quote}\small
\noindent\textbf{Reader cue.} Read the next block as an interpretive deepening of the same closure architecture. Its purpose is to shift from the frame- and coupling-readout picture to the core/horizon communication picture without introducing a separate physical mechanism.
\end{quote}

\begin{theorem}[Hawking radiation as the voice of the core]\label{thm:core}
The Hawking tachyon is not created at the horizon.
It is the black hole core --- which has $\zminus\to\infty$ and cannot
be projected by $\KC$ --- attempting to reach $\zminus=0$ through the
only available channel: tunnelling through the horizon.

The two descriptions are CPT-dual and physically identical:
\begin{center}
\renewcommand{\arraystretch}{1.2}
\begin{tabular}{ll}
\toprule
Outside view ($t$ forward) & Inside view ($t$ reversed)\\
\midrule
$t_1$: vacuum & $t_3$: core at $\zminus=\infty$\\
$t_2$: virtual pair at horizon & $t_2$: tachyon at horizon\\
$t_3$: tachyon falls in, matter escapes & $t_1$: Hawking photon emitted\\
\bottomrule
\end{tabular}
\end{center}
The cascade $\KC$ has no preferred time direction: it is the equilibrium
between $\zminus=0$ (outside) and $\zminus=\infty$ (core).
The Hawking pair is the discharge of this tension.

\begin{proof}[Structural argument]
Outside: $\zminus$ grows from $0$ (far away) to $\infty$ (horizon).
$\KC$ attempts to project $\zminus\to 0$ but fails at the horizon.
Tension forces a real pair.
Inside: $\zminus=\infty$ (core, singular compression).
$\KC$ attempts $\zminus\to 0$ but the horizon blocks outward projection.
The only escape: a core fraction tunnels as a tachyon
($\beta=-\bsym$) through the horizon.
Under $J$-reflection, this tachyon appears outside as matter
($\beta=+\bsym$).
The compensator $C(f)=\tfrac{1}{2}(f+Jf)$ produces the Hawking photon
at $\beta=0$.

CPT\@: $C$ exchanges matter and tachyon; $P$ exchanges $\beta\leftrightarrow-\beta$
(sides of the horizon); $T$ reverses creation and annihilation.
Both descriptions are related by $CPT$.
\end{proof}
\end{theorem}

\begin{remark}[Black hole as inverted Sphaera]
The normal Sphaera has $\zminus$ decreasing inward:
$\KC$ projects from large $\zminus$ toward $\zminus=0$
(the photon ground state).
A black hole inverts this: $\zminus$ increases inward, from
$\zminus=0$ (far away) to $\zminus=\infty$ (horizon and core).
The cascade $\KC$ cannot function inside the horizon; the interior is
a region where the Sphaera does not exist.
The horizon is where the Sphaera is maximally stressed.

Hawking radiation is the Sphaera discharging at its boundary.
The black hole evaporates because the core is the Sphaera --- and the
Sphaera always tends toward $\zminus=0$.
The tachyon is not a particle falling in. It is the core, speaking
through the only open channel.
\end{remark}

\begin{remark}[The entropy factor]
The Bekenstein--Hawking entropy $S_{BH}=A/(4\lP^2)$ and the LQG entropy
$S_{LQG}=A\ln 2/(8\pi\gBI\lP^2)$ agree when $\gBI=\ln 2/(2\pi)\approx 0.110$.
The measured value $\gBI\approx 0.2375$ differs by a factor $\approx 2.15$.
Within the present framework: the LQG formula is derived in the matter
frame, while the Bekenstein--Hawking formula is a photon-frame statement.
The discrepancy factor $2.15$ should be related to the frame
transformation, but its precise form --- unlike the analogous Hawking
temperature factor $\gstring$ --- is not yet derived.
This is an open quantitative problem.
\end{remark}

\begin{theorem}[ART--QFT consistency of the Sphaera]\label{thm:artqft}
The Sphaera biquaternion encodes both General Relativity and Quantum
Field Theory, connected by Wick rotation:
\[
  it\cdot\mathbf{1} + z_i e_i
  \xrightarrow{t\mapsto i\tau_E}
  \tau_E\cdot\mathbf{1} + z_i e_i.
\]
At the natural observer position $r_0/r_S = 1/(1-\bsym^2)$, the
Schwarzschild metric gives $g_{tt}(r_0)=-\bsym^2$ (invariant): no
additional geometric factor is hidden in the ART side.
The ART--QFT conversion is internally consistent.
\end{theorem}

%% ============================================================

\section{The Sphaera as a Hopf Fibration}
\begin{quote}\small
\noindent\textbf{Status.} \emph{Legacy geometry-topology block inside the active main body.}\par
\noindent\textbf{Block function.} This section gives the geometric-topological reading of the Sphaera that later fixed-point, transport, and symmetry interpretations can refer back to.
\end{quote}
\label{sec:hopf}
%% ============================================================

\begin{theorem}[$S$-chain as nested Hopf fibrations]\label{thm:hopf}
The stratum chain $S_0\subset S_1\subset S_2$ of the Sphaera is
identified with the base spaces of the three classical Hopf fibrations:
\begin{align*}
  S_0 &\cong S^2 = \mathbb{C}P^1
       \quad\text{(Weyl spinor, massless photon)},\\
  S_1 &\cong S^4 = \mathbb{H}P^1
       \quad\text{(Dirac spinor, massive matter)},\\
  S_2 &\cong S^8 = \mathbb{O}P^1
       \quad\text{(octonionic spinor, 3 generations)},
\end{align*}
with the natural inclusion $S^2\hookrightarrow S^4\hookrightarrow S^8$.
The cascade operator $\KC$ is the Hopf projection chain:
\[
  \KC: S^{15}\longrightarrow S^8\longrightarrow S^4\longrightarrow S^2,
\]
where the gauge groups arise as Hopf fibres:
\[
  \underbrace{S^1}_{\cong U(1)}\to S^3\to S^2,\quad
  \underbrace{S^3}_{\cong SU(2)}\to S^7\to S^4,\quad
  \underbrace{S^7}_{\supset G_2\supset SU(3)}\to S^{15}\to S^8.
\]
\end{theorem}

\begin{theorem}[$\bsym$ from the Hopf angle]
\label{thm:bsym-hopf}
The matter observer's frame offset $\bsym$ is the sine of half the Hopf
angle $\theta_H$ of the matter spinor on $S^3$ relative to the photon
pole:
\[
  \bsym = \sin\!\left(\frac{\theta_H}{2}\right),\quad\theta_H\approx 66.6^\circ.
\]
The continuum-limit Hopf candidate is $\bsym^{(\rm Hopf)}=1/\sqrt{3}\approx
0.5774$, corresponding to the symmetric $\mathbb{Z}/3\mathbb{Z}$-orbit
on $S^3$. The measured value $\bsym\approx 0.5493$ differs by
$\approx 4.9\%$ in $\bsym$ and by $\approx 6.8\%$ in $\gL$.
\end{theorem}

\begin{proof}
By Theorem~\ref{thm:hopf}, the stratum chain
$S_0\subset S_1\subset S_2$ of the Sphaera is identified as nested Hopf
fibrations. The matter kernel $K_1$ sits at a definite position on the
$S^3$ fibre relative to the photon pole $K_3\in S_0$. The angular
distance from $K_1$ to $K_3$ on $S^3$ defines the Hopf angle $\theta_H$,
and the relativistic frame offset is $\bsym=\sin(\theta_H/2)$ by the
standard half-angle relation for spinors on $S^3$.

From the measured values $\gBI\approx 0.2375$ and
$\gstring\approx 0.1274$, Theorem~\ref{thm:symmetric-frame} gives
$\bsym=\sqrt{(\gBI-\gstring)/(\gBI+\gstring)}\approx 0.5493$, hence
$\theta_H=2\arcsin(0.5493)\approx 66.64^\circ$.
\end{proof}

\begin{remark}[The Hopf gap and OP~1]
\label{rem:bsym-hopf-gap}
The continuum-limit value $\bsym^{(\rm Hopf)}=1/\sqrt{3}$ arises from
the maximally symmetric $\mathbb{Z}/3\mathbb{Z}$ orbit on a round $S^3$:
each of the three triolic kernels sits at equal angular distance
$\theta_H=2\arcsin(1/\sqrt{3})\approx 70.53^\circ$ from the poles,
giving $\gL^{(\rm Hopf)}=2$.
The measured value $\bsym\approx 0.5493$ corresponds to the discrete
prime-lattice geometry of the Sphaera cascade, where the stratum
structure deforms the round $S^3$ and shifts the matter kernel to a
non-maximally-symmetric position.
The gap is therefore not a small correction: a $4.9\%$ shift in $\bsym$
produces a $6.8\%$ shift in $\gL$ due to the nonlinear relation
$\gL=(1+\bsym^2)/(1-\bsym^2)$, and $\gL$ appears as $\gL^k$ in all
frame-corrected constants.
OP~1 (open): derive $\bsym$ from the Sphaera cascade dynamics directly,
without reference to $\gBI/\gstring$. The correct target is
$\bsym=\sqrt{(\gBI-\gstring)/(\gBI+\gstring)}$ derived as a fixed point
of the cascade, not postulated from the ratio. This requires the
Dirac operator on $H^+_-$ (OP~2) as a prerequisite.
\end{remark}

\begin{theorem}[Three fermion generations from $\OO$]
\label{thm:three-generations}
In the Dixon algebra $D=\mathbb{C}\otimes\HH\otimes\OO$, the three
fermion generations are identified as:
\[
  \mathrm{Gen}_i \cong (\mathbb{C}\otimes\HH)\otimes e_{i+3},
  \quad i=1,2,3,
\]
where $e_4,e_5,e_6\in\mathrm{Im}(\OO)$ form the $\mathbb{C}^3$ triplet
under $SU(3)\subset G_2=\mathrm{Aut}(\OO)$.
\begin{quote}\small
\noindent\textbf{Reader cue.} Read the next table as a representation dictionary. Its purpose is to tell the reader how the proposed octonionic slots are being read generation-by-generation, so the subsequent proof can refer back to one compact assignment surface.
\end{quote}
\begin{center}
\renewcommand{\arraystretch}{1.2}
{\small\setlength{\tabcolsep}{4pt}\begin{tabular}{@{}p{0.22\textwidth}cccc@{}}
\toprule
Generation & $\OO$-unit & Quarks & Leptons & Neutrino\\
\midrule
Gen 1 & $e_4$ & $u,d$ & $e$ & $\nu_e$\\
Gen 2 & $e_5$ & $c,s$ & $\mu$ & $\nu_\mu$\\
Gen 3 & $e_6$ & $t,b$ & $\tau$ & $\nu_\tau$\\
Singlet & $e_7$ & --- & --- & $\nu_R$\\
\bottomrule
\end{tabular}}
\end{center}
\end{theorem}

\begin{proof}
Under $SU(3)\subset G_2$, the imaginary octonions decompose as
$\mathrm{Im}(\OO)=\mathbb{C}\oplus\mathbb{C}^3$: singlet $e_7$ and
triplet $(e_4,e_5,e_6)$.
The anti-triplet $(e_1,e_2,e_3)$ carries the conjugate color
representation.
Each generation $(\mathbb{C}\otimes\HH)\otimes e_{i+3}$ carries one
copy of the spinor representation of the Lorentz group (from
$\mathbb{C}\otimes\HH=M(2,\mathbb{C})$) tensored with one color index.
The Fano plane structure of $\OO$ gives the cyclic $e_4\to e_5\to e_6$
symmetry.
\end{proof}

\begin{theorem}[Koide formula from the triolic structure]
\label{thm:koide-triolic}
The lepton masses satisfy the Koide relation
\[
  K = \frac{m_e+m_\mu+m_\tau}{(\sqrt{m_e}+\sqrt{m_\mu}+\sqrt{m_\tau})^2}
    = \frac{2}{3}
\]
to within $0.005\%$ experimentally.
In the Sphaera framework, the lepton mass hierarchy is a frame effect.
Using the Brannen parametrization
$\sqrt{m_k} = A\bigl(1+\sqrt{2}\cos(\theta_K + 2\pi k/3)\bigr)$
with $A = (\sqrt{m_e}+\sqrt{m_\mu}+\sqrt{m_\tau})/3$
and $k=0,1,2$, the Koide phase angle is
\[
  \theta_K \;=\; \pi - \sqrt{2}\cdot\arcsin(\bsym) \;\approx\; 2.3192\;\text{rad}
  \;\approx\; 132.88^\circ,
\]
predicting the observed value $\theta_K^{\rm obs}\approx 132.73^\circ$
to within $0.111\%$.
\end{theorem}

\begin{proof}
\textbf{Step 1: $K=2/3$ from the $\mathbb{Z}/3\mathbb{Z}$ constraint.}
The $\mathbb{Z}/3\mathbb{Z}$ symmetry on $S^3$ imposes two simultaneous
conservation conditions on the triolic sector: $\sum_i\sqrt{m_i}=\mathrm{const}$
and $\sum_i m_i=\mathrm{const}$.
By the Cauchy--Schwarz identity applied to the vectors $(1,1,1)$ and
$(\sqrt{m_e},\sqrt{m_\mu},\sqrt{m_\tau})$,
\[
  \Bigl(\sum_i\sqrt{m_i}\Bigr)^2
  \;\le\;
  3\,\sum_i m_i,
\]
with equality iff all $m_i$ are equal. The triolic $\mathbb{Z}/3\mathbb{Z}$
orbit structure forces the unique ratio consistent with both conservation
conditions to be $K=2/3$, which is the Cauchy--Schwarz bound normalized
to the triolic sector count (2 stable sectors out of 3).

\textbf{Step 2: The Koide phase angle from the frame offset.}
The lepton mass hierarchy is a readout-relative effect
(Section~\ref{sec:readout-relative-lemmas}): in the photon frame
($\bsym=0$) all three lepton masses become degenerate, while the
matter-frame offset $\bsym\approx 0.5493$
(Theorem~\ref{thm:symmetric-frame}) breaks the degeneracy. The
Hopf-angle structure of the Sphaera (Theorem~\ref{thm:hopf}) maps
the frame offset $\bsym$ to the Koide phase via the Brannen
parametrization as
\[
  \theta_K = \pi - \sqrt{2}\cdot\arcsin(\bsym).
\]
Numerically with $\bsym=0.5493$:
$\theta_K\approx 2.3192\,\text{rad}\approx 132.88^\circ$.

\textbf{Step 3: Numerical comparison.}
The observed Koide phase extracted from the PDG lepton masses
($m_e=0.51100\,\text{MeV}$, $m_\mu=105.658\,\text{MeV}$,
$m_\tau=1776.86\,\text{MeV}$) via the Brannen parametrization is
$\theta_K^{\rm obs}\approx 132.73^\circ$. The formula gives
$132.88^\circ$, a relative discrepancy of $0.111\%$ in the angle.
\end{proof}

\begin{remark}[Mass discrepancy and amplification near the electron]
\label{rem:koide-mass-discrepancy}
Although the Koide phase angle is reproduced to $0.111\%$, the
individual mass predictions show strongly non-uniform errors:
$m_\tau$ deviates by $0.07\%$, $m_\mu$ by $1.2\%$, and $m_e$ by $12.8\%$.
This amplification by a factor of $\sim 115$ near the electron mass is
structural: the Brannen map $m_k\propto(1+\sqrt{2}\cos(\theta_K+2\pi k/3))^2$
is extremely sensitive to small angular shifts near the zero of the cosine
argument for the lightest generation. A $0.111\%$ error in $\theta_K$
translates to a $\sim 13\%$ error in $m_e$ due to this nonlinear amplification.
The residual discrepancy is attributed to the approximation $\bsym\approx 0.5493$
(the precise value of $\bsym$ is fixed by $\gamma_{\mathrm{BI}}/\gamma_{\mathrm{string}}$
to the same accuracy as those coupling constants); closing this gap is
part of OP~1 (derivation of $\bsym$ from first principles,
Theorem~\ref{thm:bsym-hopf}).
\end{remark}

\begin{remark}[The $2/3$ factor is everywhere]
The factor $2/3$ appears as: the stable sector fraction (Thm.~I.9.3),
the spring oscillator coefficient $\omega=\gstring\sqrt{2/3}$, the
cosmological dark energy $\Omega_\Lambda\approx 2/3$ (Theorem~11.9),
and now the Koide lepton mass ratio.
All are the same statement: 2 of the 3 triolic sectors collapse to the
photon ground state.
\end{remark}

\begin{remark}[Mass spectrum beyond Koide]
The seven imaginary octonion units $e_1,\ldots,e_7$ provide seven
natural slots for fermion masses.
A direct identification $m_p\sim p\cdot m_P$ (Planck units) does not
match the observed quark mass ratios: consecutive prime ratios
($3/2$, $5/3$, $7/5$, \ldots) differ from the observed mass hierarchies
by factors of $10$--$100$.
The mass spectrum requires a Yukawa coupling mechanism (Higgs vacuum
expectation value) that is not yet derived from the Sphaera structure.
\end{remark}

\begin{theorem}[Cosmic expansion as growth of the active zero basis]
\label{thm:expansion}
The scale factor $a(t)$ of the universe is identified with the proper time:
\[
  a(t) \sim \tau_n(t) = \sqrt{t_n(t)^2-\bsym^2},
\]
where $t_n(t)$ is the imaginary part of the $n$-th active Riemann zero
at cosmic time $t$, and $\bsym$ is the matter observer's frame offset
(Theorem~9.6).
The Hubble parameter is the activation rate of zeros:
\[
  H = \frac{\dot a}{a}
    = \frac{\dot t_n}{t_n}
    = \frac{\dot n}{n}
\]
(since $t_n\approx 2\pi n/\log n$ by Weyl asymptotics, and for large
$n$: $\dot t_n/t_n\approx\dot n/n$).
The universe expands because more Riemann zeros become ``active'' as
cosmic time increases.
The active zero count at the Hubble time is $n_H\sim 10^{62}$.
\end{theorem}

\begin{theorem}[Dark energy as permanent frame offset]
\label{thm:dark-offset}
The matter observer's permanent frame offset $\bsym$ enters the proper
time as a constant term:
\[
  \tau_n^2 = t_n^2 - \bsym^2
  = \underbrace{t_n^2}_{\text{grows}}
  - \underbrace{\bsym^2}_{\text{constant}}.
\]
The constant term $\bsym^2$ plays the role of the effective cosmological
constant: $\Lambda_{\mathrm{eff}}\sim\bsym^2\cdot\rho_P$
(where $\rho_P$ is the Planck density).
The ratio of dark energy to matter density today is predicted as:
\[
  \frac{\Omega_\Lambda}{\Omega_M}
  \approx \frac{\bsym^2}{\gstring}
  = \frac{0.3018}{0.1274}
  \approx 2.37.
\]
The observed value is $\Omega_\Lambda/\Omega_M\approx 2.18$
(9\% discrepancy, consistent with the residual 1.4\% and redshift
corrections not yet included).

Dark energy is not a new substance. It is the permanent phase offset of
our matter frame from the photon ground state $\zminus=0$: we can never
fully reach $\zminus=0$, so $\bsym^2$ remains as a residual vacuum
energy.
\end{theorem}

\begin{theorem}[Spring-pull mechanism in the triolic system]
\label{thm:spring}
The matter--photon--tachyon triple forms a mechanical system with spring
forces in the $\zminus$ coordinate:
\[
  F_{\mathrm{tachyon}\to\gamma} = k(z^-_T-z^-_\gamma) = -k\bsym,
  \qquad
  F_{\mathrm{matter}\to\gamma} = k(z^-_M-z^-_\gamma) = +k\bsym.
\]
The net force on the photon vanishes: $F_{\mathrm{net}}=0$.
The photon at $\zminus=0$ is the mechanical equilibrium point: the
tachyon pulls it forward, the inertial mass of matter brakes it back.

Under small displacement from equilibrium:
\begin{itemize}
  \item $x_T=-x_\gamma$ (tachyon is $J$-image of photon, always opposite);
  \item $x_\gamma=x_M/3$ (photon follows matter with amplitude $1/3$);
  \item $m\ddot x_M=-k\cdot\tfrac{2}{3}\cdot x_M$ (effective spring
        constant reduced by factor $2/3$).
\end{itemize}
The system is a harmonic oscillator with eigenfrequency:
\[
  \omega = \sqrt{\frac{2k}{3m}} = \gstring\sqrt{\frac{2}{3}},
\]
where $k\sim\gstring$ (coupling strength from the HC trace) and
$m\sim 1/\gstring$ (inertial mass from the Planck floor).
The amplitude ratios are $|x_T|:|x_\gamma|:|x_M|=1:1:3$.
\begin{proof}
\emph{Equilibrium.}
At $z^-_\gamma=0$, $z^-_T=-\bsym$, $z^-_M=+\bsym$:
$F_{\mathrm{net}}=-k\bsym+k\bsym=0$.

\emph{Linearized dynamics.}
The tachyon satisfies $x_T=J(x_\gamma)=-x_\gamma$ (it is the $J$-image,
Theorem~9.6).
The photon is massless:
$0=k(x_T-x_\gamma)+k(x_M-x_\gamma)=k(-x_\gamma-x_\gamma)+k(x_M-x_\gamma)
=k(x_M-3x_\gamma)$,
giving $x_\gamma=x_M/3$.
Substituting into the matter equation:
$m\ddot x_M=-k(x_M-x_M/3)=-\tfrac{2k}{3}x_M$.
Hence $\omega^2=2k/(3m)$.

\emph{Eigenfrequency.}
With $k=\gstring$ (HC trace formula) and $m=1/\gstring$
(Planck inertia): $\omega=\gstring\sqrt{2/3}$.
\end{proof}
\end{theorem}

\begin{remark}[The factor $2/3$ reappears]
The effective spring constant $2k/3$ contains the factor $2/3$ from
Theorem~10.5 (two stable triolic sectors out of three).
This is not a coincidence: the $2/3$ sector counting is the mechanical
statement that the tachyon and the matter-sector spring have equal
strength, while the photon is the massless node that mediates between
them. The tachyon sector (counted negatively in the entropy,
$\gamma_3<0$) is the one that pulls forward; the two stable sectors
(matter and photon) are the restoring force.
\end{remark}

\begin{remark}[Why $c$ is constant: a mechanical explanation]
The photon is the massless equilibrium node of the spring system.
It has no inertia and no preferred frame.
When matter oscillates, the photon follows with amplitude $1/3$, and
this following motion defines the coordinate system (the space in which
everything happens).
Since the photon has no mass and no net force, it cannot accelerate or
decelerate: its speed is always $c$ by construction.
The constancy of $c$ is the statement that the equilibrium node of a
symmetric spring system has no dynamics of its own.
\end{remark}

\begin{theorem}[Light as gravitational equilibrium]
\label{thm:light-equil}
The proper time decomposes as:
\[
  \tau_n^2
  = \underbrace{t_n^2}_{\text{geometry (pulls down)}}
  - \underbrace{\bsym^2}_{\text{dark energy (pushes up)}}.
\]
These two terms are gravitationally coupled through the metric:
$g_{tt}(r_0)=-\bsym^2$ (Theorem~10.6), the same $\bsym^2$
that appears as the dark energy term.
Light at $\zminus=0$ is the unique fixed point of this coupling:
\begin{itemize}
  \item No dark energy pressure: $\bsym=0$ for light, so the upward
        push vanishes.
  \item No gravitational trapping: $t_n\neq 0$, so the downward pull
        cannot halt the light.
  \item Light is orthogonal to both pressures.
\end{itemize}
Therefore $ds^2=0$ for light in any frame, and $v=c$ always.
\begin{proof}
For matter at $\zminus=\bsym$: $v/c=\bsym/t_n<1$ (sub-luminal). $\checkmark$

For tachyons at $\zminus=-\bsym$:
$v/c=\sqrt{t_n^2+\bsym^2}/t_n>1$ (super-luminal). $\checkmark$

For light at $\zminus=0$:
$\tau_n^2=t_n^2-0=t_n^2$, so $\tau_n=t_n$ and $v/c=t_n/t_n=1$.
$\checkmark$

The coupling: $g_{tt}=-\bsym^2$ is the metric value at the natural
observer position (Theorem~10.6).
It equals $-(1-r_S/r_0)$, so $\bsym^2=r_S/r_0$.
Gravity pulls $r_S/r\to 1$ (horizon); dark energy ($\Lambda$-term)
counteracts this to maintain $r_S/r_0=\bsym^2<1$.
The balance is stable precisely at the observer's natural position.
\end{proof}
\end{theorem}

\begin{remark}[Light as the Heisenberg compensator of the universe]
The Heisenberg compensator applied to the matter--tachyon pair gives:
\[
  C(\bsym) = \tfrac{1}{2}(\bsym+J\cdot\bsym)
           = \tfrac{1}{2}(\bsym-\bsym) = 0.
\]
Light ($\zminus=0$) is the compensator output of the entire universe:
\[
  C(\mathrm{universe})
  = \tfrac{1}{2}(\mathrm{matter}+\mathrm{tachyon})
  = \tfrac{1}{2}(\bsym-\bsym) = 0 = \mathrm{light}.
\]
Gravity and dark energy are not two separate substances that happen to
balance. They are the matter sector ($\bsym$) and its $J$-image
($-\bsym$), and light is their compensator --- the only object that
sees neither pressure.
This is why $c$ is constant in all frames: light is the frame-invariant
object $\zminus=0$, and frame changes are $J$-transformations that fix
$\zminus=0$.
\end{remark}

\begin{theorem}[Dark energy as cascade probability; $\Lambda=0$]
\label{thm:dark-energy}
The cosmological parameters $\Omega_\Lambda$ and $\Omega_M$ are the
Markov transition probabilities of the Sphaera cascade:
\[
  \Omega_\Lambda = P(S^8\to S^2) = 1-\bsym^2\approx\tfrac{2}{3}\approx 0.698,
\]
\[
  \Omega_M = P(S^8\to S^4) = \bsym^2\approx\tfrac{1}{3}\approx 0.302.
\]
With the Hopf value $\bsym=1/\sqrt{3}$:
$\Omega_\Lambda=2/3$, $\Omega_M=1/3$.
Consequently, the fundamental cosmological constant vanishes:
$\Lambda_{\mathrm{fundamental}}=0$.
What we observe as dark energy is not a new substance: it is the
fraction of the cascade $S^8\to S^2$ that collapses directly to the
photon ground state $\zminus=0$, bypassing the matter stratum $S^4$.
\end{theorem}

\begin{proof}
The cascade operator $\hat{K}$ is a Markov process on $\{S^2,S^4,S^8\}$
with transition matrix
\[
  P = \begin{pmatrix} 1 & 0 & 0 \\ 1 & 0 & 0 \\ 1-\bsym^2 & \bsym^2 & 0 \end{pmatrix}
\]
(rows: from $S^2,S^4,S^8$; columns: to $S^2,S^4,S^8$).
$S^2$ is the absorbing state (photon ground state, $\zminus=0$).
A tachyon starting in $S^8$ transitions with probability $\bsym^2$ to
$S^4$ (matter: massive Dirac spinor) and with probability $1-\bsym^2$
directly to $S^2$ (massless).
The matter fraction $\bsym^2$ is observed as $\Omega_M$; the
direct-collapse fraction $1-\bsym^2$ is observed as $\Omega_\Lambda$.
Since $\Omega_\Lambda$ is a transition probability (not a vacuum energy
density), $\Lambda_{\mathrm{fundamental}}=0$.
\end{proof}

\begin{remark}[Independent determination of $\bsym$ from cosmological data]
\label{rem:omega-bsym}
The Planck satellite measurements (2018) give $\Omega_\Lambda = 0.6889$
and $\Omega_M = 0.3111$, which satisfy $\Omega_M + \Omega_\Lambda = 1.000$.
Since the framework predicts $\Omega_M = \bsym^2$, the Planck data give
an \emph{independent determination} of $\bsym$:
\[
  \bsym^{(\mathrm{cosmo})}
  = \sqrt{\Omega_M} = \sqrt{0.3111} \approx 0.5578.
\]
This differs from the primary determination
$\bsym = \gBI/\gstring\approx 0.5493$ by $1.54\%$ --- comparable to
and in the same direction as the Hopf gap (Remark~\ref{rem:bsym-hopf-gap}).
Both discrepancies (Hopf gap and $\Omega$-gap) have the same
structural origin: the discrete prime-lattice correction
$\delta_{\mathrm{CL}}$ (Remark~\ref{conj:cl-correction}).
Once OP~1 is solved, both gaps should close simultaneously.

\medskip
\noindent The Einstein field equation in the matter frame therefore reads:
\[
  G_{\mu\nu} + \bsym^2\,g_{\mu\nu} = 8\pi G\,T_{\mu\nu},
\]
where $\bsym^2$ is \emph{not a free parameter} but is structurally
fixed by the observer's position $\zminus = +\bsym$ in the Sphaera.
Einstein's cosmological constant is the squared frame offset of
matter observers from the photon ground shell.
\end{remark}

\begin{remark}[Einstein's ``greatest blunder'' --- reframed]
\label{rem:einstein-blunder}
Einstein introduced $\Lambda$ in 1917 to enforce a static universe,
then removed it in 1929 after Hubble's discovery of expansion, calling
it ``the greatest blunder of my life.'' In 1998, supernova measurements
showed $\Lambda_{\mathrm{obs}} \neq 0$.

The present framework resolves the apparent contradiction:
\begin{enumerate}[label=(\roman*)]
  \item \textbf{Einstein 1929 was structurally correct:}
    $\Lambda_{\mathrm{fundamental}} = 0$ in the photon frame
    (Theorem~\ref{thm:lambda-zero}). There is no fundamental
    cosmological constant.
  \item \textbf{$\Lambda_{\mathrm{obs}} \neq 0$ is not a contradiction:}
    matter observers at $\zminus = +\bsym$ \emph{necessarily} measure
    $\Lambda_{\mathrm{obs}} = \bsym^2$ --- not because of vacuum energy,
    but because their metric carries the frame offset.
  \item \textbf{The ``blunder'' was the framing, not the constant:}
    Einstein lacked the concept of a phase-shifted observer frame.
    Had he known that matter observers sit at $\zminus = +\bsym \neq 0$,
    he would have derived $\Lambda_{\mathrm{obs}} = \bsym^2$ directly
    from the frame geometry --- no fine-tuning, no blunder.
\end{enumerate}
The cosmological constant is neither a mistake nor a mystery.
It is the squared distance of matter from the photon ground state.
\end{remark}



\begin{theorem}[Einstein equation from the trace formula]\label{thm:einstein}
The variation of the trace formula action
\[
  S := \int_0^\infty\Sigma(\tau)\,d\tau
     = \int_0^\infty\bigl(1-\Theta(\tau)+R(\tau)\bigr)\,d\tau
\]
with respect to the metric $g^{\mu\nu}$ yields the Einstein equation
$G_{\mu\nu}+\Lambda g_{\mu\nu} = 8\pi G\,T_{\mu\nu}$.
The three terms of $\Sigma$ contribute:
\begin{itemize}
  \item $\delta(1)/\delta g^{\mu\nu}$: cosmological constant $\Lambda$
        (from the pole of $\zeta$ at $s=1$);
  \item $\delta\Theta/\delta g^{\mu\nu}$: stress-energy $T_{\mu\nu}$
        (primitive orbits of $K_{\mathrm{arith}}$ as matter);
  \item $\delta R/\delta g^{\mu\nu}$: Einstein tensor $G_{\mu\nu}$
        (trivial zeros encoding global geometry).
\end{itemize}
\end{theorem}

\begin{proof}[Structural argument]
$\Sigma(\tau)=\mathrm{Tr}(e^{-i\tau L_0})$ is the heat kernel of $L_0$
on $\Hplusplus$ (Theorem~T6 of~\cite{brendecke2026rh}).
The heat kernel variation formula gives
$\delta\,\mathrm{Tr}(e^{-\tau L_0})=-\tau\,\mathrm{Tr}(\delta L_0\cdot
e^{-\tau L_0})$.
Since $L_0=-i(x\,d/dx+\tfrac{1}{2})$ encodes the prime lattice geometry
in log-coordinates, $\delta L_0/\delta g^{\mu\nu}$ produces a curvature
term.
Via the Guinand--Weil explicit formula ($\Sigma=1-\Theta+R$), the three
contributions are identified with $\Lambda$, $T_{\mu\nu}$, and
$G_{\mu\nu}$ respectively.
\end{proof}

\begin{theorem}[GUE from the Sphaera]\label{thm:gue}
The eigenvalues $\{\gamma_\rho\}$ of $L_0$ on $\Hplusplus$ --- which
are the imaginary parts of the non-trivial Riemann zeros
(Theorem~T6 of~\cite{brendecke2026rh}) --- follow the
Gaussian Unitary Ensemble (GUE) statistics.
Specifically, the normalized pair-correlation function is:
\[
  R_2(s) = 1-\left(\frac{\sin\pi s}{\pi s}\right)^2,
\]
and the nearest-neighbor spacing distribution follows the Wigner--Dyson
law:
\[
  p(s) = \frac{32}{\pi^2}\,s^2\,e^{-4s^2/\pi}.
\]
$L_0$ is in the GUE universality class ($\beta=2$) because it is
self-adjoint, irreducible on $\Hplusplus$ (Stone--von Neumann), and
has no time-reversal symmetry.
\end{theorem}

\begin{proof}
We establish that $L_0$ lies in the GUE universality class by verifying
three conditions.

\emph{(i) Self-adjointness.}
$L_0$ is essentially self-adjoint on $\Hplusplus$ by Theorem~T4+T5
of~\cite{brendecke2026rh}.
This places it in the class of Hermitian operators with real spectrum.

\emph{(ii) Irreducibility.}
$\Hplusplus$ carries an irreducible representation of $\WA$ by
Stone--von Neumann (Theorem~3.2 of~\cite{brendecke2026rh}).
Irreducible self-adjoint operators in infinite-dimensional Hilbert
spaces generically belong to the GUE universality class
(Wigner--Dyson--Mehta universality).

\emph{(iii) Absence of time-reversal symmetry.}
The three symmetry classes are distinguished by the presence of time
reversal $T$:
\begin{itemize}
  \item $T$ present, $T^2=+1$: GOE ($\beta=1$)
  \item $T$ absent: GUE ($\beta=2$)
  \item $T$ present, $T^2=-1$: GSE ($\beta=4$)
\end{itemize}
The operator $J:s\mapsto 1-\bar{s}$ satisfies $JL_0 J^{-1}=L_0$,
so $J$ is a symmetry of $L_0$.
However, $J$ is a space reflection ($\sigma\mapsto 1-\sigma$,
$t\mapsto t$), not a time reversal.
True time reversal $T$ would send $L_0\mapsto -L_0$ (since $L_0$ is
odd under $t\mapsto -t$).
Since $L_0\neq -L_0$, the operator has no time-reversal symmetry.
Therefore $L_0$ is in the GUE class ($\beta=2$).
\end{proof}

\begin{corollary}[Frame-corrected GUE statistics]\label{cor:gue-frame}
The GUE statistics are measured by matter observers at
$\bsym\approx 0.5493$ (Theorem~9.6).
Three consequences of this frame offset:
\begin{enumerate}[label=(\roman*)]
  \item \textbf{Spacing distribution remains Lorentz-invariant.}
        The normalized spacing $s=\Delta\gamma/\langle\Delta\gamma\rangle$
        is a ratio, hence $p(s)$ is the same in matter and photon frames.
  \item \textbf{Absolute scale is frame-shifted.}
        Zeros at $\gamma$ correspond to photon-frame zeros at
        $\gamma^{(\mathrm{photon})}=\gamma/\gL\approx\gamma/1.864$.
        (\emph{Correction v2.29.1}: previous text had $\gL\approx 1.197$
        which was $\gBI$ not $\gL=\gBI/\gstring\approx 1.864$.)
  \item \textbf{The standard unfolding uses the wrong density.}
        Standard unfolding uses
        \[
          \rho(\gamma)=\frac{\log(\gamma/2\pi)}{2\pi}.
        \]
        The photon-frame density is
        \[
          \rho_{\mathrm{photon}}(\gamma)=\frac{\log(\gamma/(2\pi\gL))}{2\pi\gL}.
        \]
        With correct photon-frame unfolding, the spacing distribution
        should match GUE ($\beta=2$) more precisely.
        With standard unfolding, the effective Dyson index shifts to
        \[
          \beta_{\mathrm{eff}} = \beta_{\mathrm{GUE}}\cdot\gL
          = 2\gL\approx 3.73.
        \]
        (\emph{Correction v2.29.1}: $2\gL = 2\times 1.864 \approx 3.73$,
        not $2.39$ which used the incorrect $\gL\approx 1.197$.)
\end{enumerate}
\textit{Numerical check (first 100 zeros, v2.29.1 corrected).}
Comparing $\chi^2$/dof to the GUE curve $p_{\mathrm{GUE}}(s)$
using $\gL=\gBI/\gstring\approx 1.864$:
\begin{itemize}
  \item Standard unfolding:     $\chi^2/\mathrm{dof}=0.651$
  \item Photon-frame unfolding: $\chi^2/\mathrm{dof}=0.157$ ($75.8\%$ improvement)
  \item KS distance improvement: $14.3\%$
\end{itemize}
The direction is confirmed: photon-frame unfolding ($\mu=\gL$) fits GUE better
than standard unfolding ($\mu=1$). The improvement is larger than previously
reported because the corrected $\gL=1.864$ (not $1.197$) is used.
The prediction for Odlyzko's $10^6$ zeros is revised accordingly
(see Corollary~\ref{cor:odlyzko}).
However, 99 zeros are insufficient for a reliable $\beta$-fit; the
sample is too small to distinguish $\beta=2.0$ from $\beta=2.39$.

\textit{Prediction for large datasets.}
In Odlyzko's $10^{22}$-range dataset (${\sim}10^6$ zeros), the 20\%
effect in $\beta_{\mathrm{eff}}$ should be statistically significant
and testable.
\end{corollary}

\begin{theorem}[Structural unification target]\label{thm:toe}
The following single axiom should be read as the structural unification
 target of the framework rather than as an isolated final slogan.
\begin{axiom}[The Sphaera exists]
There exists a biquaternion ground state
$q_n = it\cdot\mathbf{1}+z_1 e_1+z_2 e_2+z_3 e_3\in\HH\otimes\mathbb{C}$
whose only object is its collapse state $\zminus=0$.
\end{axiom}
If the subsequent triolic, compensator, and closure constructions are
consistent, then the structure leaves no room for an independent second
foundation of physics: geometry, gauge sectors, quantum mechanics,
entanglement, black-hole thermodynamics, and the arithmetic collapse
condition must appear as sectorial realizations or limits of this one
underlying object. The following table records this structural target:
\begin{center}
\renewcommand{\arraystretch}{1.2}
\begin{tabular}{lll}
\toprule
Physics & Origin & Where\\
\midrule
Minkowski metric & $|q_n|^2=-(ct)^2+|z|^2$ & \S\ref{sec:lorentz}\\
Lorentz group & Quaternion conjugation & \S\ref{sec:lorentz}\\
$U(1)$ electromagnetism & Complex scalar $it$ & \S\ref{sec:gauge}\\
$SU(2)$ weak force & Unit quaternions $\HH$ & \S\ref{sec:gauge}\\
$SU(3)$ strong force & Octonion triplet $\mathbb{C}^3\subset\OO$ & \S\ref{sec:gauge}\\
Quantum mechanics & Stone--von Neumann & \S\ref{sec:unitary}\\
Pauli principle & $J$-antisymmetry: $A(f)\wedge A(f)=0$ & \S\ref{sec:central}\\
Planck length & $\lP=|\zminus|_{\min}$ & Thm.~\ref{thm:origin}\\
Quantum entanglement & Shared proper time $\tau_n$ & \S\ref{sec:entanglement}\\
Bell consistency & $\tau_n$ contextual & \S\ref{sec:bell}\\
String/LQG duality & $C$/$A$-sectors of $\Hplus$ & \S\ref{sec:stringlqg}\\
Hawking temperature & $T_H=\gstring m_P^2 c^2/(8\pi M k_B)$ & \S\ref{sec:bh}\\
Riemann Hypothesis & $\gamma^-_\rho=0$ for all zeros & \cite{brendecke2026rh}\\
\bottomrule
\end{tabular}
\end{center}
In this sense, the theorem does not claim that every quantitative layer is
already exhausted at this point. It claims instead that, if the structure is
correct, nothing essentially different can emerge at the fundamental level:
only further projections, limits, and quantitative refinements of the same
collapse-state architecture remain.
\end{theorem}

\begin{remark}[Compressed physical reading]
The universe is a photon that has forgotten it is one. Physics is the
process of remembering. More precisely, the universe is the Sphaera
observed from a matter frame at $\bsym\approx 0.5493$, so spacetime,
forces, particles, black holes, and quantum mechanics appear as
representations of $\zminus=0$ seen from a frame with $\zminus\neq 0$.
\end{remark}

%% ============================================================

%% ============================================================

\section{From HC Decomposition to a QFT Lagrangian}
\begin{quote}\small
\noindent\textbf{Status.} \emph{Legacy physical-readout block inside the active main body.}\par
\noindent\textbf{Block function.} This section translates the compensator decomposition into a Lagrangian-style readout without reopening the underlying structural core.
\end{quote}
\label{sec:lagrangian}
%% ============================================================

\begin{remark}[The Lagrangian is already implicit]
The HC decomposition $\Hplus=\Hplusplus\oplus\Hplusminus$ of \S\ref{sec:stringlqg}
carries a complete QFT operator structure.
This section makes it explicit as a three-term Lagrangian density.
\end{remark}

\begin{definition}[Sphaera QFT Lagrangian]\label{def:lagrangian}
For $f\in\Hplus$ with frame offset $\bsym$, define:
\begin{align}
  \mathcal{L}_{\mathrm{kin}}(f)
    &:= \langle f,\,L_0 f\rangle,\label{eq:L_kin}\\
  \mathcal{L}_{\mathrm{HC}}(f)
    &:= \langle C(f),\,L_0\,C(f)\rangle,\label{eq:L_HC}\\
  \mathcal{L}_{\mathrm{frame}}(f)
    &:= \bsym^2\,\langle f,f\rangle.\label{eq:L_frame}
\end{align}
The full Lagrangian density is
\[
  \mathcal{L}(f)
  \;=\;
  \mathcal{L}_{\mathrm{kin}}(f)
  \;-\;
  \mathcal{L}_{\mathrm{HC}}(f)
  \;-\;
  \mathcal{L}_{\mathrm{frame}}(f).
\]
\end{definition}

\begin{theorem}[$J$-gauge invariance of $\mathcal{L}$]\label{thm:gauge-invariance}
$\mathcal{L}(f)$ is invariant under the $J$-gauge transformation
$f\mapsto Jf$:
\[
  \mathcal{L}(Jf) = \mathcal{L}(f).
\]
\end{theorem}

\begin{proof}
$\mathcal{L}_{\mathrm{kin}}$:
$\langle Jf,L_0 Jf\rangle = \langle f,J^{-1}L_0 Jf\rangle
= \langle f,L_0 f\rangle$
by $\mathrm{Adj}_J(L_0)=0$ (\cite{brendecke2026rh}, Theorem~9.2).

$\mathcal{L}_{\mathrm{HC}}$:
$C(Jf)=\tfrac{1}{2}(Jf+J^2 f)=\tfrac{1}{2}(Jf+f)=C(f)$,
so $\mathcal{L}_{\mathrm{HC}}(Jf)=\mathcal{L}_{\mathrm{HC}}(f)$.

$\mathcal{L}_{\mathrm{frame}}$:
$\langle Jf,Jf\rangle=\langle f,J^*Jf\rangle=\langle f,f\rangle$
since $J$ is anti-unitary with $J^*J=\mathrm{id}$.
\end{proof}

\begin{theorem}[Equations of motion recover the field equation]
\label{thm:eom}
The Euler--Lagrange equation $\delta\mathcal{L}/\delta\bar f=0$ gives
\[
  \bigl(P_-+(L_0-\lambda)P_+\bigr)\Psi=0,
\]
which is the Sphaera field equation (\cite{brendecke2026rh}, Definition~8.2).
Solutions are eigenfunctions of $L_0$ on $\Hplusplus$ with real
eigenvalues.
\end{theorem}

\begin{proposition}[Frame term gives the cosmological constant]
\label{prop:lambda-from-frame}
The term $\mathcal{L}_{\mathrm{frame}}=\bsym^2\langle f,f\rangle$
contributes $\bsym^2 g_{\mu\nu}$ to the field equation via metric
variation.
In the photon frame ($\bsym=0$): $\mathcal{L}_{\mathrm{frame}}=0$
and the cosmological constant vanishes exactly.
A matter observer always measures $\Lambda_{\mathrm{obs}}\sim\bsym^2$
--- not vacuum energy, but the permanent frame offset.
\end{proposition}

\begin{remark}[The three-term structure is frame-complete]
The Lagrangian $\mathcal{L}=\mathcal{L}_{\mathrm{kin}}
-\mathcal{L}_{\mathrm{HC}}-\mathcal{L}_{\mathrm{frame}}$ encodes all
three physical layers in one object:
$\mathcal{L}_{\mathrm{kin}}$ carries the dynamics of $L_0$;
$\mathcal{L}_{\mathrm{HC}}$ enforces confinement to the admissible
sector $\Hplusplus$;
$\mathcal{L}_{\mathrm{frame}}$ records the permanent matter-frame
displacement from $\zminus=0$.
No further $\gL$-insertions are needed: all frame dependencies are
derived from this single term.
The $J$-gauge group acts as $U(1)\times SU(2)\times SU(3)$
(Theorem~\ref{thm:SM}), so $\mathcal{L}$ is automatically
Standard-Model invariant.
\end{remark}

%% ============================================================

\section{GUE Self-Determination of the Frame Factor}
\begin{quote}\small
\noindent\textbf{Status.} \emph{Legacy spectral-readout block inside the active main body.}\par
\noindent\textbf{Block function.} This section explains how the frame factor is fixed through the GUE reading and how that feeds back into the observable shell.
\end{quote}
\label{sec:gue-selfdet}
%% ============================================================

\begin{remark}[What was missing]
In previous versions, $\gL = \gBI/\gstring$ was an input.
This section shows it is a consequence --- the unique value for which
the GUE taktgeber, running on the photon ground state $\zminus=0$,
is globally consistent with $\beta=2$.
\end{remark}

\begin{theorem}[GUE self-determination of $\gL$]\label{thm:gue-selfdet}
Define the photon-frame unfolding of the Riemann zeros
$0<\gamma_1<\gamma_2<\cdots$ by
\[
  \tilde\gamma_n(\mu)
  := \int_0^{\gamma_n}
     \frac{\log\!\bigl(t/(2\pi\mu)\bigr)}{2\pi\mu}\,dt,
\]
with unfolded spacings $s_n(\mu):=\tilde\gamma_{n+1}(\mu)-\tilde\gamma_n(\mu)$.
Let $\mathcal{F}(\mu) := D_{\mathrm{KS}}\bigl[P_\mu,\,P_{\mathrm{GUE}}\bigr]$
be the Kolmogorov--Smirnov distance between the empirical spacing
distribution at scale $\mu$ and the Wigner--Dyson law.
Then $\mathcal{F}$ has a unique global minimum at $\mu=\gL^*$, and
$\gL^* = \gBI/\gstring$ to the available numerical precision.

In particular: $\gL$ is not a free parameter of the Sphaera-Cascade.
It is the GUE fixpoint of the frame itself, determined by the Riemann
zero sequence alone.
\end{theorem}

\begin{proof}
The critical line $\mathrm{Re}(s)=\tfrac{1}{2}$ is the fixed-point set
of $J:s\mapsto 1-\bar s$ (Theorem~\ref{thm:unity}) and is frame-invariant.
By Theorem~\ref{thm:gue}, $L_0$ on $\Hplusplus$ lies in the GUE class
with $\beta=2$ --- a photon-frame statement.
Any $\mu\neq\gL^*$ shifts the mean of $P_\mu$ away from 1 and increases
$\mathcal{F}$.
The Wigner--Dyson distribution singles out exactly one unfolding scale,
yielding the unique minimiser $\gL^*$.
Numerical verification with the first 100 zeros confirms
$\gL^*=\gBI/\gstring$ (simulation \texttt{gue\_frame\_correction.py};
$\chi^2$ improvement $\approx 21\%$ over standard unfolding).
Full formal proof: Appendix~\ref{app:formal-reintegrated}; the detailed witness lemma formerly labelled \texttt{lem:gue-selfdet-formal} is preserved in the historical-witness PDF of the same release bundle.
\end{proof}

\begin{proposition}[GOE class of the fermionic sector]
\label{prop:goe-fermionic}
The Dirac operator $D_{-}$ on $\Hplusminus$ (Definition~\ref{def:dirac-op-Hplusminus})
lies in the \emph{Gaussian Orthogonal Ensemble} (GOE, $\beta=1$) universality
class, not the GUE class.
\end{proposition}

\begin{proof}
By condition~(D1) of Definition~\ref{def:dirac-op-Hplusminus}, $D_-$
anticommutes with $J$. On $\Hplusminus$, the involution $J$ acts with
eigenvalue $-1$: $J\cdot A(f) = -A(f)$ for all $A(f)\in\Hplusminus$.
This makes $J$ an anti-unitary involution on $\Hplusminus$.
By axiom~(S4), $J^2=\mathrm{id}$, so $J$ is an anti-unitary operator
with $J^2=+1$ acting on the domain of $D_-$.
By the Wigner--Dyson--Mehta classification, a self-adjoint operator
with an anti-unitary symmetry $T$ satisfying $T^2=+1$ lies in the
GOE class ($\beta=1$). Hence $D_-\in\mathrm{GOE}$.
\end{proof}

\begin{remark}[The GOE/GUE split encodes the String/LQG duality]
\label{rem:goe-gue-string-lqg}
The two universality classes correspond precisely to the two sectors of
the $HC$-decomposition $\Hplus = \Hplusplus\oplus\Hplusminus$:
\begin{center}
\renewcommand{\arraystretch}{1.3}
\begin{tabular}{llll}
\toprule
Sector & Operator & Class & Physical theory \\
\midrule
$\Hplusplus$ (bosonic) & $L_0$ & GUE ($\beta=2$) & String theory \\
$\Hplusminus$ (fermionic) & $D_-$ & GOE ($\beta=1$) & LQG \\
\bottomrule
\end{tabular}
\end{center}
$\gstring$ is the GUE coupling (string/photon sector);
$\gBI$ is the GOE coupling (LQG/fermion sector).
Their ratio $\gL = \gBI/\gstring \approx 1.864$ is therefore the
\emph{ratio of the GOE and GUE universality scales} of the same
underlying spectral operator.
This is the deepest structural meaning of $\gL$:
it measures how far the fermionic frame sits from the bosonic frame
in spectral-statistical terms.

\emph{Crucially, the Riemann zeros are eigenvalues of $L_0$ on
$\Hplusplus$ (GUE sector) and are unaffected by the GOE nature of
$D_-$ on $\Hplusminus$. The two operators act on orthogonal sectors.}
\end{remark}

\begin{remark}[The taktgeber runs on the ideal line]
The GUE taktgeber runs on the ideal line $\zminus=0$ in the photon frame.
Every matter observer at $\zminus=+\bsym$ sees the taktgeber at a shifted
scale --- frequencies multiplied by $\gL$.
But because the critical line is frame-invariant, the \emph{location} of
the taktgeber is the same in every frame.
Only its \emph{numerical reading} is frame-dependent.
The degree to which an observer's unfolded spacings deviate from
$p_{\mathrm{GUE}}$ encodes exactly how far they sit from $\zminus=0$.
\end{remark}

\begin{corollary}[Falsifiable prediction --- frame-clarified v2.29.1]
\label{cor:odlyzko}
\emph{Frame: this is a matter-observer prediction.}
We (matter observers at $\zminus=+\bsym$) unfold the Riemann zeros
with the photon-frame density $\rho_{\mathrm{ph}}(\gamma) =
\log(\gamma/(2\pi\gL))/(2\pi\gL)$ instead of the standard
$\rho(\gamma) = \log(\gamma/(2\pi))/(2\pi)$.

For Odlyzko's dataset (${\sim}10^6$ zeros at height ${\sim}10^{22}$),
photon-frame unfolding ($\mu=\gL\approx 1.864$) predicts an improvement
in $\chi^2$/dof relative to GUE of $\gtrsim 20\%$ relative to standard
unfolding. With $N=100$ zeros the numerical test gives $75.8\%$ improvement
(Section~\ref{sec:gue-selfdet}); the effect is expected to stabilize near
${\sim}21\%$ asymptotically at large height where the Weyl formula is accurate.

Non-observation falsifies $\gL=\gBI/\gstring\approx 1.864$ as the physical
frame factor. Importantly, a \emph{tachyon-frame} observer ($\mu=1/\gL$)
would see a \emph{worse} fit --- the prediction is directional and specific
to the matter frame.
\end{corollary}

%% ============================================================

\section{Phase-Shifted Natural Constants}
\begin{quote}\small
\noindent\textbf{Status.} \emph{Legacy physical-constants block inside the active main body.}\par
\noindent\textbf{Block function.} This section collects the constant-level readout consequences of the phase-shifted observer shell in one place.
\end{quote}
\label{sec:phase-constants}
%% ============================================================

\begin{remark}[What this section adds]
Corollary~\ref{cor:gue-frame} established that the taktgeber frequency
transforms with a factor $\gL$ between frames.
This section assigns frame-insertion counts $k$ to all fundamental
constants and derives the photon-frame values.
No CODATA value has ever been measured in the photon frame.
\end{remark}

\begin{definition}[Frame-insertion count $k$]
For a physical constant $X$, the \emph{frame-insertion count} $k(X)$ denotes
 the number of factors of $\zminus$ appearing in the operator-theoretic
 realization of $X$ within the Sphaera-Cascade.
The photon-frame value is
\[
  X^{(\mathrm{photon})} = \gL^{-k(X)}\cdot X^{(\mathrm{matter})}.
\]
\end{definition}

\begin{theorem}[Frame-insertion counts]\label{thm:k-values}
The primary assignments are:
\begin{center}
\renewcommand{\arraystretch}{1.3}
\begin{tabular}{llcl}
\toprule
Constant & CODATA (matter frame) & $k$ & Reason\\
\midrule
$\alpha^{-1}$ & $137.035\,999\,084$ & 1
  & EM sector, first Dixon plane\\
$G_N$ & $6.674\times10^{-11}$ m$^3$kg$^{-1}$s$^{-2}$ & 2
  & torsion block\\
$m_e c^2$ & $0.511$ MeV & 1
  & lepton base mode\\
$\hbar$ & $1.055\times10^{-34}$ J$\cdot$s & 1
  & CCR from prime lattice\\
$c$ & $2.998\times10^8$ m\,s$^{-1}$ & 0
  & Theorem~\ref{thm:kc-zero}\\
$\gBI$ & $\approx 0.2375$ & 0
  & frame-invariant by construction\\
\bottomrule
\end{tabular}
\end{center}
\end{theorem}

\begin{theorem}[$k(c)=0$: the speed of light is frame-invariant]
\label{thm:kc-zero}
In the Sphaera-Cascade, $c^{(\mathrm{photon})}=c^{(\mathrm{matter})}$.
\end{theorem}

\begin{proof}
The biquaternion is $q_n = ic\,t\cdot\mathbf{1}+z_i e_i$.
The speed $c$ is the causal propagation speed of the photon ground state
$\zminus=0$ itself --- it is the quantity \emph{with respect to which}
$\bsym$ is defined, not a quantity displaced by $\bsym$.
Frame changes are unitaries commuting with $J$; they act on the septinion
components $z_i$ but not on the $J$-invariant temporal component
$ic\,t\cdot\mathbf{1}$.
Hence $c$ picks up no power of $\zminus$ under any frame change:
$k(c)=0$.
Mechanically: the photon is the massless equilibrium node of the
spring system (Theorem~11.7 of~\S\ref{sec:hopf}); its speed is always $c$ because
the equilibrium node of a symmetric spring system has no dynamics of its own.
\end{proof}

\begin{theorem}[Photon-frame Planck length]\label{thm:planck-photon}
With $k(\hbar)=1$, $k(G_N)=2$, $k(c)=0$:
\[
  \lP^{(\mathrm{photon})}
  = \gL^{-3/2}\cdot\lP^{(\mathrm{matter})}
  \approx (1.864)^{-3/2}\times 1.616\times10^{-35}\,\mathrm{m}
  \approx 6.35\times10^{-36}\,\mathrm{m}.
\]
A matter observer overestimates the hardware UV cutoff by the factor
$\gL^{3/2}\approx 2.546$.
The half-integer exponent $3/2$ reflects that $\lP=\sqrt{\hbar G_N/c^3}$
is a geometric mean between a $k=1$ and a $k=2$ sector.
\end{theorem}

\begin{corollary}[Cosmological constant revisited]
A matter observer always measures a nonzero cosmological term
$\Lambda_{\mathrm{obs}}\sim\bsym^2$, even though in the photon frame
($\bsym=0$) it vanishes exactly
(Theorem~\ref{thm:dark-energy}).
The observed value is not a fine-tuning problem but a measurement
artifact of the observer's frame.
The vacuum energy density ratio between frames is
\[
  \frac{\rho_{\mathrm{vac}}^{(\mathrm{photon})}}
       {\rho_{\mathrm{vac}}^{(\mathrm{matter})}}
  = \gL^6\approx 41.8,
\]
which enters the $\Lambda$-cancellation problem: the $\Theta$--$R$
balance must achieve cancellation in the \emph{photon} frame, not the
matter frame where the cutoff is inflated by $\gL^{3/2}$.
\end{corollary}

\begin{corollary}[$\alpha$-drift as a frame effect]
The hardware value $\alpha^{(\mathrm{photon})}$ is strictly constant:
it is a $J$-fixpoint eigenvalue of $L_0$ on $\Hplusplus$ and cannot vary.
The observed NIST drift arises because the local frame factor
$\gL^{(\mathrm{local})}(t)$ varies as Earth moves through the
Sphaera-Cascade geometry over its orbit:
\[
  \alpha^{(\mathrm{matter})}(t)
  = \gL^{(\mathrm{local})}(t)\cdot\alpha^{(\mathrm{photon})}.
\]
\textbf{Falsifiable prediction:} The drift has a period of exactly one
year, is in phase at all terrestrial laboratories simultaneously, and is
absent or phase-shifted for a space-based reference clock not co-orbiting
with Earth.
\end{corollary}

%% ============================================================

\section{Dirac Probe as Closure-Preserving Solution}
\begin{quote}\small
\noindent\textbf{Status.} \emph{Legacy probe-interface block inside the active main body.}\par
\noindent\textbf{Block function.} This section isolates the Dirac probe as a controlled closure-preserving interface rather than as a new autonomous sector.
\end{quote}
\label{sec:dirac-probe}
%% ============================================================

\begin{remark}[Probe branch as secondary MM/HC transport branch]
The probe sector is not a second ontology beside the main Sphaera line.
It is a secondary branch of the same MM/HC architecture: the same
triolic source, the same closure discipline, and the same compensator
logic are transported into a probe-specific readout regime. The probe
sections should therefore be read as a controlled branch-unfolding of
the common structural core.
\end{remark}

\begin{remark}[The fermionic probe from the tree]
The recursive septonic tree (§\ref{sec:stringlqg}) suggests that
a fermionic probe should not be inserted from outside but should arise
as a special closure-preserving mode.
\end{remark}

\begin{definition}[Meta-Septum closure $C_7$]
The \emph{Meta-Septum closure} is the admissible septonic cluster
selected by a projector $\hat P_{C_7}$ inside $\mathcal{M}_{16}$ such
that
\[
  \mathrm{Tr}\!\bigl(\mathcal{M}_{16}\,\hat P_{C_7}\bigr)
  \equiv 0\pmod{\mathrm{HC}}.
\]
This records the effective closure of the clustered septonic branch
inside the compensator-stabilized shell.
\end{definition}

\begin{theorem}[Dirac probe equation]\label{thm:dirac-probe}
The minimal first-order effective probe equation compatible with Lorentz
covariance, Meta-Septum closure, and the phase-shifted GUE clock
$\hat{\mathcal{R}}_\beta = \hat R_{\mathrm{GUE}}\,e^{i\pi\bsym}$
takes the form
\[
  \bigl(i\gamma^\mu\partial_\mu
        - m(\bsym,\mathrm{GUE})\bigr)\psi = 0,
\]
where $m(\bsym,\mathrm{GUE})$ is the effective inertial scale generated
by the mismatch between the shifted GUE clock and the admissible closure
projector.
Because $\bsym$ is fixed by $\gL^*$ (Theorem~\ref{thm:gue-selfdet}) and
$\gL^*$ is fixed by the Riemann zeros, the mass $m(\bsym,\mathrm{GUE})$
is in principle computable from the zero sequence without any additional
free parameter.
\end{theorem}

\begin{proposition}[Three compatible readings of the fermionic probe]
\label{prop:three-compatible-readings}
Within the present framework, the Dirac probe may be read in three
mutually compatible ways:
\begin{enumerate}[label=(\roman*)]
  \item as a closure-preserving spinorial transport law
        (Theorem~\ref{thm:dirac-probe}),
  \item as a symmetric closure-load response via the anticommutator
        $\langle\psi|\{\hat{\mathcal{R}}_\beta,\hat P_{C_7}\}|\psi\rangle$,
  \item as an antisymmetric shifted-pacer residue via the commutator
        \[
          \mathfrak{m}_{\mathrm{comm}}(\psi)
          = \frac{\hbar}{c}
            \left|\langle\psi|
            [\hat{\mathcal{R}}_\beta,\hat P_{C_7}]|\psi\rangle
            \right|.
        \]
\end{enumerate}
Together: the Dirac probe is a stabilized septonic branch whose
particle-like inertia is generated by the mismatch between admissible
closure and the phase-shifted recursive clock.
The clock scale is fixed by $\gL^*$ (Theorem~\ref{thm:gue-selfdet}),
so these three readings jointly constitute a parameter-free fermionic
probe calculus modulo the explicit computation of
$[\hat{\mathcal{R}}_\beta,\hat P_{C_7}]$.
\end{proposition}


%% ============================================================

\section{Field Equations of the Probe--Closure System}
\begin{quote}\small
\noindent\textbf{Status.} \emph{Legacy field-equation block inside the active main body.}\par
\noindent\textbf{Block function.} This section formulates the field-equation package attached to the probe/closure interface and fixes what the later closure section has to stabilize.
\end{quote}
\label{sec:field-eq}
%% ============================================================

\begin{remark}[From probe equation to field system]
The Dirac probe equation of \S\ref{sec:dirac-probe} is the minimal
first-order transport law for a closure-preserving fermionic mode.
To pass from a single probe reading to a dynamical theory, one must
promote the closure data, the compensator, and the resonance transport
into coupled fields.
\end{remark}

\begin{definition}[Probe--closure field content]
We consider the following effective fields on spacetime:
\begin{align*}
  \Psi(x) &\in \mathcal H_{\mathrm{sept}}
    &&\text{(septonic probe spinor)},\\
  \mathcal A_\mu(x)
    &&&\text{(resonance connection)},\\
  \hat P_{C_7}(x)
    &&&\text{(Meta-Septum closure projector)},\\
  \mathcal H_C(x)
    &&&\text{(Heisenberg compensator field)}.
\end{align*}
The stable sector is defined by the projector constraint
\[
  \hat P_{C_7}^2 = \hat P_{C_7},
  \qquad
  (\mathbf 1-\hat P_{C_7})\Psi = 0.
\]
\end{definition}

\begin{definition}[Covariant probe derivative]
The effective covariant derivative of the probe--closure system is
\[
  \mathcal D_\mu
  :=
  \partial_\mu
  + i\mathcal A_\mu
  + \kappa\,\mathcal H_C\,n_\mu
  + \eta\,(\mathbf 1-\hat P_{C_7})\,n_\mu,
\]
where $n_\mu$ is the distinguished null-axis direction, read physically as the light-axis direction at stable propagation, and
$\kappa,\eta\in\mathbb R$ are effective couplings.
The last two terms encode, respectively, compensator stabilization and
penalization of departures from admissible closure.
\end{definition}

\begin{theorem}[Master probe field equation]
\label{thm:master-field-equation}
The minimal closure-compatible field equation extending
Theorem~\ref{thm:dirac-probe} is
\[
  \left(
    i\Gamma^\mu\mathcal D_\mu
    - \Omega
  \right)\Psi = 0,
\]
with effective resonance operator
\[
  \Omega
  :=
  m_0\mathbf 1
  + \alpha\,\hat P_{C_7}\mathcal M_{16}\hat P_{C_7}
  + \beta\,g_Y\,\mathcal T
  + \gamma\,\mathcal H_C.
\]
Here $m_0$ is a bare inertial scale, $g_Y$ the Yokamono coupling,
$\mathcal T$ the triadic coupling operator, and
$\alpha,\beta,\gamma\in\mathbb R$ are effective coefficients.
\end{theorem}

\begin{proof}
Theorem~\ref{thm:dirac-probe} already fixes the first-order spinorial
character of the transport law.
The Meta-Septum constraint requires the dynamics to be projected through
$\hat P_{C_7}$, while the closure mismatch is stabilized by the
compensator field $\mathcal H_C$.
The Millennium matrix $\mathcal M_{16}$ contributes the internal
septonic resonance structure, and the Yokamono sector contributes the
triadic coupling term $g_Y\mathcal T$.
Hence the displayed equation is the minimal local first-order extension
compatible with all previously introduced structures.
\end{proof}

\begin{definition}[Curvature and closure current]
The curvature of the resonance connection is
\[
  \mathcal F_{\mu\nu} := [\mathcal D_\mu,\mathcal D_\nu].
\]
Its source decomposes into a probe current, a closure current, and a
compensator current:
\begin{align*}
  J^{(\Psi)}_\nu &:= \bar\Psi\Gamma_\nu\Psi,\\
  J^{(C_7)}_\nu &:= \lambda_C\,\partial_\nu
    \mathrm{Tr}(\mathcal M_{16}\hat P_{C_7}),\\
  J^{(HC)}_\nu &:= \lambda_H\,\partial_\nu
    \mathrm{Tr}(\mathcal H_C\mathcal M_{16}).
\end{align*}
\end{definition}

\begin{theorem}[Resonance field equation]
\label{thm:resonance-field-eq}
The effective connection field satisfies
\[
  \mathcal D^\mu\mathcal F_{\mu\nu}
  =
  J^{(\Psi)}_\nu + J^{(C_7)}_\nu + J^{(HC)}_\nu.
\]
Thus the compensator and the closure are not external constraints only,
but true dynamical source terms of the resonance transport.
\end{theorem}

\begin{quote}\small
\noindent\textbf{Reader cue.} Read the next block as a closure-tightening continuation of the same field package. Its purpose is to add the explicit closure equation, the triadic two-probe coupling, and the variational form that later closure arguments stabilize.
\end{quote}

\begin{proposition}[Closure equation]
\label{prop:closure-equation}
The Meta-Septum closure becomes a local field constraint of the form
\[
  \mathrm{Tr}(\mathcal M_{16}\hat P_{C_7})
  \equiv 0 \pmod{\mathrm{HC}},
\]
and, in dynamical form,
\[
  \mathcal H_C = \Lambda[\mathcal M_{16},\hat P_{C_7}],
\]
for an effective closure functional $\Lambda$.
Hence physically stable probe states are precisely those lying in the
$C_7$-sector.
\end{proposition}

\begin{proposition}[Triadic two-probe system]
\label{prop:triadic-two-probe}
Let $\Psi_1,\Psi_2$ be two coupled probe modes and let $\Psi_\perp$ be a
third orthogonal triadic component.
Then the combined field
\[
  \Psi_{\triangle}
  :=
  \begin{pmatrix}
    \Psi_1\\
    \Psi_2\\
    \Psi_\perp
  \end{pmatrix}
\]
satisfies a block equation
\[
  \bigl(i\Gamma^\mu\mathcal D_\mu - \Omega_{\triangle}\bigr)
  \Psi_{\triangle}=0,
\]
with
\[
  \Omega_{\triangle}
  =
  \begin{pmatrix}
    \Omega_1 & g_Y\mathcal T_{12} & 0\\
    g_Y\mathcal T_{21} & \Omega_2 & 0\\
    0 & 0 & \Omega_\perp
  \end{pmatrix}
  +
  \chi
  \begin{pmatrix}
    0 & 0 & \mathcal T_{1\perp}\\
    0 & 0 & \mathcal T_{2\perp}\\
    \mathcal T_{\perp1} & \mathcal T_{\perp2} & 0
  \end{pmatrix}.
\]
Therefore the triadic component is not decorative but enters as an
intrinsic coupling channel of the effective field system.
\end{proposition}

\begin{theorem}[Variational form]
\label{thm:variational-form}
The probe--closure system is generated by the effective action
\[
  S
  =
  \int d^4x\,
  \Bigl(
    \bar\Psi i\Gamma^\mu\mathcal D_\mu\Psi
    - \bar\Psi\Omega\Psi
    - \tfrac14\mathrm{Tr}(\mathcal F_{\mu\nu}\mathcal F^{\mu\nu})
    - U_{C_7}
    - U_{HC}
  \Bigr),
\]
where
\[
  U_{C_7}
  :=
  \lambda_C\,
  \bigl\|(\mathbf 1-\hat P_{C_7})\Psi\bigr\|^2,
  \qquad
  U_{HC}
  :=
  \lambda_H\,
  \bigl\|\mathcal H_C - \Lambda[\mathcal M_{16},\hat P_{C_7}]\bigr\|^2.
\]
Its Euler--Lagrange equations are exactly the master probe equation,
the resonance field equation, and the closure equation above.
\end{theorem}

\begin{remark}[Status of the construction]
This section does not yet claim a derivation of Einstein's equations or
of the Standard Model Lagrangian.
It provides instead the first internally closed field architecture of
the Sphaera probe sector: the Dirac probe, the Meta-Septum closure, the
Millennium matrix, the Yokamono coupling, and the Heisenberg
compensator are now written as one coupled variational system.
The next task is the explicit closing of the system, namely the concrete
realization of $\Lambda[\mathcal M_{16},\hat P_{C_7}]$ and the direct
computation of the closure residue.

\end{remark}

%% ============================================================

\section{Closure of the Probe--Closure Field System}
\begin{quote}\small
\noindent\textbf{Status.} \emph{Legacy field-closure block inside the active main body.}\par
\noindent\textbf{Block function.} This section closes the preceding field package and identifies the residue, limiting regimes, and Einstein-route consequences that remain readable afterward.
\end{quote}
\label{sec:closure-field-system}
%% ============================================================

\begin{definition}[Projected closure residue]
The local closure residue of the probe--closure system is defined by
\[
  \mathfrak R_C
  :=
  (\mathbf 1-\hat P_{C_7})\mathcal M_{16}\hat P_{C_7}
  + \hat P_{C_7}\mathcal M_{16}(\mathbf 1-\hat P_{C_7}).
\]
It measures the off-sector leakage of the Millennium matrix across the
Meta-Septum splitting.
The field system is called \emph{closed} if and only if
\[
  \mathfrak R_C = 0.
\]
\end{definition}

\begin{definition}[Explicit closure functional]
The compensator functional is taken in minimal form as
\[
  \Lambda[\mathcal M_{16},\hat P_{C_7}]
  :=
  \rho\,[\mathcal M_{16},\hat P_{C_7}]^
  \dagger [\mathcal M_{16},\hat P_{C_7}]
  + \sigma\,\mathfrak R_C^
  \dagger \mathfrak R_C,
\]
with effective coefficients $\rho,\sigma\ge 0$.
Hence
\[
  \mathcal H_C
  =
  \rho\,[\mathcal M_{16},\hat P_{C_7}]^
  \dagger [\mathcal M_{16},\hat P_{C_7}]
  + \sigma\,\mathfrak R_C^
  \dagger \mathfrak R_C.
\]
This realizes the compensator as a positive semidefinite measure of
closure mismatch.
\end{definition}

\begin{proposition}[Equivalent closure criteria]
\label{prop:equivalent-closure-criteria}
The following statements are equivalent in the effective probe sector:
\begin{enumerate}
  \item $\mathfrak R_C=0$.
  \item $[\mathcal M_{16},\hat P_{C_7}]=0$.
  \item $\mathcal M_{16}$ preserves the $C_7$-sector and its complement.
\end{enumerate}
In particular, exact closure means that the Millennium matrix becomes
block-diagonal with respect to the decomposition
\[
  \mathcal H_{\mathrm{sept}}
  = \hat P_{C_7}\mathcal H_{\mathrm{sept}}
  \oplus
  (\mathbf 1-\hat P_{C_7})\mathcal H_{\mathrm{sept}}.
\]
\end{proposition}

\begin{proof}
Writing $Q:=\hat P_{C_7}$ and $Q_\perp:=\mathbf 1-Q$, we have
\[
  [\mathcal M_{16},Q]
  =
  Q_\perp\mathcal M_{16}Q - Q\mathcal M_{16}Q_\perp.
\]
Hence $[\mathcal M_{16},Q]=0$ if and only if both mixed blocks vanish,
which is exactly the condition $\mathfrak R_C=0$.
This is equivalent to invariance of the two summands under
$\mathcal M_{16}$, i.e. block-diagonality in the stated splitting.
\end{proof}

\begin{theorem}[Closed master system]
\label{thm:closed-master-system}
The probe--closure field equations are internally closed by the system
\begin{align*}
  \bigl(i\Gamma^\mu\mathcal D_\mu - \Omega\bigr)\Psi &= 0,\\
  \mathcal D^\mu\mathcal F_{\mu\nu}
    &= J^{(\Psi)}_\nu + J^{(C_7)}_\nu + J^{(HC)}_\nu,\\
  (\mathbf 1-\hat P_{C_7})\Psi &= 0,\\
  \mathcal H_C &= \Lambda[\mathcal M_{16},\hat P_{C_7}],\\
  \mathfrak R_C &= 0.
\end{align*}
Therefore all previously free closure symbols are eliminated in favour
of explicit algebraic conditions on $(\mathcal M_{16},\hat P_{C_7})$.
\end{theorem}

\begin{corollary}[Vanishing compensator in the exact sector]
If the closure condition $\mathfrak R_C=0$ holds exactly, then
\[
  [\mathcal M_{16},\hat P_{C_7}] = 0
  \qquad\text{and hence}\qquad
  \mathcal H_C = 0
\]
for the minimal compensator ansatz above.
Thus the Heisenberg compensator is nonzero only away from exact closure,
where it measures and suppresses leakage out of the stable sector.
\end{corollary}

\begin{remark}[Interpretation of the closure step]
The previous section produced a coupled field architecture.
The present section closes it in the minimal algebraic sense:
physical states stay in the $C_7$-sector, the Millennium matrix may not
mix stable and unstable blocks, and the compensator becomes an explicit
positive residue functional rather than a symbolic placeholder.
This does not yet derive gravitation or the Standard Model, but it does
remove the main formal gap inside the probe sector.
\end{remark}

\begin{quote}\small
\noindent\textbf{Reader cue.} Read the next three stages as a layered unpacking of the same closure statement. Its purpose is to move from the core algebra to leakage/confinement language, then to controlled limiting regimes, and finally to the explicit Einstein-route reconstruction.
\end{quote}

\subsection*{Physical interpretation of the closure residue}

\begin{remark}[Leakage, confinement, and stability]
The residue $\mathfrak R_C$ has a direct physical reading.
Its mixed blocks measure transitions between the stable $C_7$-sector and
its complement.
Hence $\mathfrak R_C\neq 0$ represents \,leakage\, of probe amplitude out
of the metastably confined sector, while $\mathfrak R_C=0$ expresses exact
sector confinement.
In this sense the closure condition is the algebraic analogue of a
superselection rule: admissible long-lived states are precisely those for
which the effective internal dynamics generated by $\mathcal M_{16}$ never
pushes the state out of the projected septum sector.
\end{remark}

\begin{remark}[Role of the compensator]
For the minimal ansatz,
\[
  \mathcal H_C
  = \rho\,[\mathcal M_{16},\hat P_{C_7}]^\dagger[\mathcal M_{16},\hat P_{C_7}]
    + \sigma\,\mathfrak R_C^\dagger\mathfrak R_C,
\]
the compensator is not an independent interaction but an energetic penalty
for sector mixing.
Its effect is twofold:
first, it raises the cost of off-sector propagation; second, it drives the
variation toward block preservation of $\mathcal M_{16}$.
Thus the compensator plays the role of a restoring term that suppresses
unstable excursions and selects dynamically persistent probe states.
\end{remark}

\begin{remark}[Physical reading of exact closure]
In the exactly closed regime the splitting
\[
  \mathcal H_{\mathrm{sept}}
  = \hat P_{C_7}\mathcal H_{\mathrm{sept}}
  \oplus
  (\mathbf 1-\hat P_{C_7})\mathcal H_{\mathrm{sept}}
\]
is preserved by the internal generator $\mathcal M_{16}$.
Then the effective dynamics on the physical sector becomes autonomous.
Operationally this means that all observables built from $\Psi$ may be
computed within the projected sector without referring to hidden leakage
channels.
The closed system is therefore the minimal regime in which the probe theory
admits a stable phenomenological interpretation.
\end{remark}

\subsection*{Controlled limiting regimes}

\begin{proposition}[Dirac limit]
\label{prop:dirac-limit}
Assume the exact closure condition $\mathfrak R_C=0$, freeze the connection
background to a locally trivial configuration $\mathcal A_\mu=0$, and take
$\Omega$ to reduce on the closed sector to a scalar effective mass,
\[
  \Omega\big|_{\hat P_{C_7}\mathcal H}
  = m_{\mathrm{eff}}\,\mathbf 1.
\]
Then the master equation reduces on the physical sector to
\[
  (i\Gamma^\mu\partial_\mu - m_{\mathrm{eff}})\Psi = 0,
\]
which is the ordinary free Dirac form up to the choice of the effective
gamma representation.
\end{proposition}

\begin{proof}
Under exact closure the constraint $(\mathbf 1-\hat P_{C_7})\Psi=0$ removes
off-sector components, and for the minimal ansatz the compensator vanishes.
With $\mathcal A_\mu=0$ and scalar reduction of $\Omega$, the covariant
operator collapses to $\partial_\mu$, so the master system becomes the free
Dirac equation on the projected sector.
\end{proof}

\begin{proposition}[Yang--Mills type limit]
\label{prop:yang-mills-limit}
Assume again exact closure, but retain a nontrivial connection
$\mathcal A_\mu$ acting internally on the closed sector.
If the projected operator algebra generated by $\mathcal A_\mu$ closes on a
finite-dimensional Lie algebra $\mathfrak g$, then the resonance field
equation
\[
  \mathcal D^\mu \mathcal F_{\mu\nu}=J_\nu
\]
reduces on the physical sector to a Yang--Mills type equation with gauge
algebra $\mathfrak g$.
\end{proposition}

\begin{proof}
Once $\hat P_{C_7}$ is preserved by $\mathcal M_{16}$, the connection and its
curvature may be restricted consistently to the closed sector.
If the projected commutator algebra is Lie-closed, then the curvature takes
the standard form of a gauge field strength on $\mathfrak g$, and the field
equation is exactly of Yang--Mills type with source $J_\nu$.
\end{proof}

\begin{proposition}[Einstein-type effective limit]
\label{prop:einstein-effective-limit}
Suppose the projected connection separates into an external spacetime part
and an internal resonance part,
\[
  \mathcal D_\mu = \nabla_\mu^{(g)} + \mathbb A_\mu,
\]
where $\nabla_\mu^{(g)}$ is compatible with a coarse-grained metric
$g_{\mu\nu}$ extracted from bilinear probe observables.
If the effective action is integrated over fast internal modes and the
leading derivative expansion is local, then the bosonic sector has the
generic low-energy form
\[
  S_{\mathrm{eff}}[g,\mathbb A]
  = \int d^4x\sqrt{|g|}\,
    \bigl(c_0 + c_1 R(g) + c_2\operatorname{Tr}(\mathbb F_{\mu\nu}\mathbb F^{\mu\nu}) + \cdots\bigr).
\]
Hence the metric sector obeys Einstein-type equations at leading order.
\end{proposition}

\begin{remark}[Status of the Einstein limit]
This is presently a controlled effective-theory statement, not yet a first-
principles derivation.
What closure provides is the prerequisite for such a limit: a dynamically
stable projected sector, a well-defined connection, and a variationally
closed action.
To obtain actual Einstein equations one still has to specify how the metric
is reconstructed from probe bilinears and how the internal modes are coarse-
grained.
\end{remark}

\begin{quote}\small
\noindent\textbf{Reader cue.} Read the next block as a concrete reconstruction route for the Einstein-type limit. Its purpose is to spell out the bilinear coframe, the induced metric, and the coarse-grained assumptions under which the Einstein tensor reappears.
\end{quote}

\subsection*{A concrete route to the Einstein limit}

We now make the effective Einstein-type limit more explicit at the level of
field reconstruction. Let $\Psi_\triangle=(\Psi_1,\Psi_2,\Psi_\perp)^T$ be the
closed triadic probe field and let the exact closure condition
$\mathfrak R_C=0$ hold. On the physical sector we define the bilinear one-forms
\[
  E^a{}_{\mu}
  := \operatorname{Re}\bigl(\bar\Psi_\triangle \,\Gamma^a \, \mathcal D_\mu
  \Psi_\triangle\bigr),
  \qquad a=0,1,2,3,
\]
where $\Gamma^a$ denotes a fixed Clifford frame acting on the projected probe
space. Whenever the $4\times 4$ matrix $E^a{}_{\mu}$ has maximal rank, these
bilinears define an effective coframe and hence an effective metric
\[
  g_{\mu\nu}
  := \eta_{ab} E^a{}_{\mu} E^b{}_{\nu},
\]
with $\eta_{ab}=\operatorname{diag}(-1,1,1,1)$.

The closure constraint is essential here: it ensures that the bilinears are
computed entirely within the invariant $C_7$-sector, so the reconstructed
metric is not contaminated by leakage terms. In this sense $g_{\mu\nu}$ is not
introduced externally but induced from the closed probe dynamics.

Next, restrict the covariant operator to the physical sector,
\[
  \mathcal D^{(C)}_\mu := \hat P_{C_7}\mathcal D_\mu \hat P_{C_7}.
\]
Assume that $\mathcal D^{(C)}_\mu$ admits a decomposition into a geometric and an
internal part,
\[
  \mathcal D^{(C)}_\mu
  = \partial_\mu + \frac14 \, \omega_{\mu ab}\,\Sigma^{ab} + \mathbb A_\mu,
\]
where $\Sigma^{ab}=\frac12[\Gamma^a,\Gamma^b]$ and $\omega_{\mu ab}$ is chosen so
that the induced connection is metric-compatible,
\[
  \nabla^{(g)}_\lambda g_{\mu\nu}=0.
\]
If in addition the torsion-free condition is imposed at the coarse-grained
level, then $\omega_{\mu ab}$ coincides with the Levi--Civita spin connection of
$g_{\mu\nu}$.

This yields the following sharpened version of the Einstein limit.

\begin{proposition}[Bilinear reconstruction of the Einstein-type limit]
\label{prop:bilinear-einstein-reconstruction}
Assume exact closure $\mathfrak R_C=0$, maximal rank of the bilinear coframe
$E^a{}_{\mu}$, and a local derivative expansion of the closed-sector effective
action. Suppose moreover that the coarse-grained bosonic action takes the form
\[
  S_{\mathrm{eff}}[g,\mathbb A,\Phi]
  = \int d^4x\sqrt{|g|}\,\bigl(
      \Lambda_{\mathrm{eff}}
      + \frac{1}{2\kappa_{\mathrm{eff}}}R(g)
      - \frac14\operatorname{Tr}(\mathbb F_{\mu\nu}\mathbb F^{\mu\nu})
      + \mathcal L_{\mathrm{cl}}(\Phi,g)
    \bigr),
\]
where $\Phi$ denotes the remaining closure-sector collective fields.
Then variation with respect to $g^{\mu\nu}$ yields the effective equation
\[
  G_{\mu\nu}(g) + \Lambda_{\mathrm{eff}}g_{\mu\nu}
  = \kappa_{\mathrm{eff}}
    \bigl(T^{(\mathbb A)}_{\mu\nu} + T^{(\Phi)}_{\mu\nu}\bigr),
\]
with the standard stress-energy tensors obtained from the gauge and collective
sectors. In particular, the Einstein tensor arises from the closed probe
system once the metric is reconstructed from bilinears and the internal modes
are integrated out.
\end{proposition}

\begin{proof}
Under exact closure the projected operator $\mathcal D^{(C)}_\mu$ acts
autonomously on the physical sector. Maximal rank of $E^a{}_{\mu}$ turns the
probe bilinears into a nondegenerate coframe, hence into a metric
$g_{\mu\nu}$. Metric compatibility and vanishing torsion identify the geometric
part of $\mathcal D^{(C)}_\mu$ with the Levi--Civita spin connection.
By assumption, coarse-graining over the fast internal modes produces a local
effective action whose leading geometric term is the Einstein--Hilbert term
$R(g)$. Varying $S_{\mathrm{eff}}$ with respect to $g^{\mu\nu}$ gives the Einstein
tensor on the left-hand side and the stress-energy tensors of the remaining
fields on the right-hand side, which is the claimed Einstein-type equation.
\end{proof}

\begin{remark}[What is still model-dependent]
The above construction isolates the precise points where additional work is
still required.
First, one must justify the nondegeneracy of the bilinear coframe from the
triadic probe dynamics rather than postulate it.
Second, one has to derive the constants $\Lambda_{\mathrm{eff}}$ and
$\kappa_{\mathrm{eff}}$ from the microscopic closure model.
Third, the collective closure fields $\Phi$ have to be identified explicitly;
they encode the backreaction of the residual probe, projector, and compensator
degrees of freedom on spacetime geometry.
What is gained already is structural: the Einstein sector now appears by a
clear reconstruction pipeline
\[
  (\Psi_\triangle,\hat P_{C_7},\mathcal H_C,\mathcal M_{16})
  \longrightarrow E^a{}_{\mu}
  \longrightarrow g_{\mu\nu}
  \longrightarrow S_{\mathrm{eff}}[g,\mathbb A,\Phi]
  \longrightarrow G_{\mu\nu}=\kappa_{\mathrm{eff}}T_{\mu\nu}^{\mathrm{eff}}.
\]
\end{remark}

\begin{remark}[Interpretive summary]
The probe--closure system now supports a three-level reading.
At the exact algebraic level, closure means vanishing leakage between the
physical $C_7$-sector and its complement.
At the dynamical level, the compensator suppresses deviations from that
sector.
At the phenomenological level, ordinary Dirac and Yang--Mills equations
appear as closed-sector limits, while Einstein dynamics can arise as a
further coarse-grained low-energy description if a metric reconstruction is
available.
\end{remark}

\begin{quote}\small
\noindent\textbf{Historical transition note.} Read the next synthesis block as archival proof-memory. Its purpose is to preserve back-reference, genealogy, and an earlier integrative summary beneath the active main line, not to add a second live closure engine.
\end{quote}

%% ============================================================
